% !TeX program = xelatex
% !TeX root = main.tex
%
% 2026 MathorCup 大数据竞赛 A 题论文
% 啤酒酿造化合物分析与啤酒风味预测
%
% 编译（在 latex/ 目录下执行）：
%   latexmk -xelatex main.tex          % 推荐，自动处理 bibtex 与多次编译
% 或手动：
%   xelatex main.tex
%   bibtex main
%   xelatex main.tex
%   xelatex main.tex
%
% 说明：
%   1. 编译需要 TeX Live（含 ctex 与 fandol 字体）；本仓库不附带 TeX 安装。
%   2. 图片统一放在 ../img/，由 src/ 下 Python 脚本生成；
%      图片不存在时 \figauto 会自动显示占位框，不影响编译。
%   3. 正文中所有 [TODO: ...] 与“待定”均为待填内容，定稿前须清零。

\documentclass[UTF8,a4paper,12pt,fontset=fandol]{ctexart}
% fontset=fandol 随 TeX Live 安装，Overleaf 亦可直接用；
% 若本机已装 Noto CJK，可改为 fontset=ubuntu。

% ============================================================
% preamble.tex
% 2026 MathorCup 大数据竞赛 A 题论文导言区
% 编译引擎：xelatex（中文由 ctexart / xeCJK 提供）
% 说明：本文件只放宏包、全局格式与常用宏，不放正文。
% ============================================================

% ---------- 页面与版式 ----------
\usepackage{geometry}
\geometry{left=2.6cm,right=2.6cm,top=2.5cm,bottom=2.5cm}
\linespread{1.3}
\usepackage{fancyhdr}
\pagestyle{fancy}
\fancyhf{}
\fancyfoot[C]{\thepage}
\renewcommand{\headrulewidth}{0pt}
\setlength{\emergencystretch}{2em}

% ---------- 数学 ----------
\usepackage{amsmath,amssymb}
\usepackage{bm}
\allowdisplaybreaks[3]
\numberwithin{equation}{section}
\numberwithin{figure}{section}
\numberwithin{table}{section}
\DeclareMathOperator*{\argmin}{arg\,min}
\DeclareMathOperator*{\argmax}{arg\,max}
\newcommand{\R}{\mathbb{R}}
\newcommand{\E}{\mathbb{E}}
\newcommand{\Var}{\operatorname{Var}}
\newcommand{\Cov}{\operatorname{Cov}}
\newcommand{\Corr}{\operatorname{Corr}}
\newcommand{\norm}[1]{\left\lVert #1 \right\rVert}
\newcommand{\abs}[1]{\left\lvert #1 \right\rvert}
\newcommand{\inner}[2]{\left\langle #1,#2 \right\rangle}

% ---------- 一套符号：向量、矩阵、希腊字母 ----------
% 约定：粗体小写表示向量，粗体大写表示矩阵。
\newcommand{\vx}{\bm{x}}
\newcommand{\vy}{\bm{y}}
\newcommand{\bz}{\bm{z}}
\newcommand{\bw}{\bm{w}}
\newcommand{\vb}{\bm{b}}
\newcommand{\bs}{\bm{s}}
\newcommand{\bzero}{\bm{0}}
\newcommand{\bX}{\bm{X}}
\newcommand{\bY}{\bm{Y}}
\newcommand{\bZ}{\bm{Z}}
\newcommand{\bW}{\bm{W}}
\newcommand{\bB}{\bm{B}}
\newcommand{\bA}{\bm{A}}
\newcommand{\bR}{\bm{R}}
\newcommand{\bM}{\bm{M}}
\newcommand{\bI}{\bm{I}}
\newcommand{\bD}{\bm{D}}
\newcommand{\bL}{\bm{L}}
\newcommand{\bU}{\bm{U}}
\newcommand{\bV}{\bm{V}}
\newcommand{\bP}{\bm{P}}
\newcommand{\bT}{\bm{T}}
\newcommand{\bbeta}{\bm{\beta}}
\newcommand{\btheta}{\bm{\theta}}
\newcommand{\bmu}{\bm{\mu}}
\newcommand{\bSigma}{\bm{\Sigma}}
\newcommand{\bepsilon}{\bm{\epsilon}}
\newcommand{\vxt}{\tilde{\bm{x}}}
\newcommand{\yhat}{\hat{y}}
\newcommand{\Dcal}{\mathcal{D}}
\newcommand{\Lcal}{\mathcal{L}}
\newcommand{\Ncal}{\mathcal{N}}
\newcommand{\Jcal}{\mathcal{J}}
\newcommand{\Tcal}{\mathcal{T}}
\newcommand{\Mcal}{\mathcal{M}}

% ---------- 图 ----------
\usepackage{graphicx}
% 图片统一放在项目根目录的 img/ 下，由 src/ 的 Python 脚本生成。
% 同时保留 img/，便于从项目根目录编译时也能找到图。
\graphicspath{{../img/}{img/}}
\DeclareGraphicsExtensions{.pdf,.png,.jpg,.jpeg}

% ---------- 表 ----------
\usepackage{array}
\usepackage{booktabs}
\usepackage{tabularx}
\usepackage{longtable}
\renewcommand{\arraystretch}{1.18}

% ---------- 算法 ----------
\usepackage{xcolor}
\usepackage[ruled,vlined,linesnumbered]{algorithm2e}
\SetAlgorithmName{算法}{算法}{算法列表}
\SetKwInOut{Input}{输入}
\SetKwInOut{Output}{输出}

% ---------- 代码 ----------
\usepackage{listings}
\renewcommand{\lstlistingname}{代码}
\definecolor{codeblue}{RGB}{25,70,160}
\definecolor{codegreen}{RGB}{0,110,60}
\definecolor{codered}{RGB}{170,30,30}
\definecolor{codegray}{RGB}{110,110,110}
\lstdefinestyle{pycode}{
  language=Python,
  basicstyle=\ttfamily\small,
  keywordstyle=\color{codeblue}\bfseries,
  commentstyle=\color{codegreen},
  stringstyle=\color{codered},
  numbers=left,
  numberstyle=\tiny\color{codegray},
  numbersep=6pt,
  frame=single,
  breaklines=true,
  showstringspaces=false,
  columns=fullflexible,
  keepspaces=true,
  tabsize=4,
  captionpos=b
}
\lstdefinestyle{shellcode}{
  language=bash,
  basicstyle=\ttfamily\small,
  keywordstyle=\color{codeblue}\bfseries,
  commentstyle=\color{codegreen},
  stringstyle=\color{codered},
  numbers=left,
  numberstyle=\tiny\color{codegray},
  numbersep=6pt,
  frame=single,
  breaklines=true,
  showstringspaces=false,
  columns=fullflexible,
  keepspaces=true,
  tabsize=4,
  captionpos=b
}
\lstset{style=pycode}
% 注意：代码块内请使用英文注释，避免 CJK 字符在 listings 中断行异常。

% ---------- 参考文献与链接 ----------
\usepackage[numbers,sort&compress]{natbib}
\usepackage[hidelinks,bookmarksnumbered=true]{hyperref}

% ---------- 常用占位宏 ----------
% 定稿前全文搜索 TODO 与“待定”，逐项清零。
\newcommand{\todo}[1]{\textcolor{red!75!black}{\textbf{[TODO: #1]}}}
\newcommand{\tbd}{\textcolor{red!75!black}{待定}}

% \figauto{文件名}{图题}{标签}
% 用法：\figauto{fig_q1_framework.pdf}{问题一框架}{fig:q1framework}
% 若 ../img/ 或 img/ 下存在该文件则自动插图；否则显示带文件名的占位框。
% 这样骨架在图片尚未生成时也能编译，图片生成后无需改正文。
\newcommand{\figauto}[4][0.82\textwidth]{%
  \begin{figure}[htbp]
    \centering
    \IfFileExists{../img/#2}{%
      \includegraphics[width=#1]{#2}%
    }{%
      \IfFileExists{img/#2}{%
        \includegraphics[width=#1]{#2}%
      }{%
        \fbox{\begin{minipage}[c][4.2cm][c]{0.78\textwidth}
          \centering\small
          图待生成：\texttt{\detokenize{#2}}\\
          由 \texttt{src/} 下脚本输出到 \texttt{img/}
        \end{minipage}}%
      }%
    }%
    \caption{#3}
    \label{#4}
  \end{figure}%
}

% ---------- 章节标题格式 ----------
\ctexset{
  section = {format = \Large\bfseries},
  subsection = {format = \large\bfseries},
  subsubsection = {format = \normalsize\bfseries},
  paragraph = {format = \normalsize\bfseries}
}


\begin{document}

\title{\bfseries 啤酒酿造化合物分析与啤酒风味预测\\[0.35em]
  \large ——基于多源化合物数据的分级评分、关键因素识别与风格条件生成}
\author{参赛队号：\underline{\hspace{2.5cm}}\quad 队员：\underline{\hspace{4.5cm}}}
\date{}
\maketitle
\thispagestyle{empty}

% ============================================================
% 摘要
% 只写方法与结论的骨架；所有数字与具体结论用 \todo 占位。
% ============================================================

\begin{center}
  {\zihao{4}\bfseries 摘\quad 要}
\end{center}
\vspace{0.4em}

本文围绕商用啤酒风味代谢数据集，研究化合物谱与感官评分之间的映射关系、关键风味化合物的识别，以及给定风格下的化合物谱生成与跨风格迁移问题。针对附件数据样本量有限、化合物维度高、缺失值与异常值并存、感官评分为 0--6 离散等级等特点，本文采用“基线—改进—最终”的递进式建模路线，并在每一问中同时报告预测精度、稳定性与可解释性；同时引入题目建议的 Kaggle Top Beer Information 数据集，在 22 个风格层级做外部效度校验（覆盖 5,558 行中的 85.6\%），并明确其无法逐酒合并的限制。

针对问题一，本文以 \texttt{4.csv} 的 50 个感官维度为监督信号，把 \texttt{1.csv} 的 231 个化合物特征映射到连续感官潜变量，再校准为 0--6 分级。基线为逐维均值/岭回归；改进模型加入缺失指示、分位数变换与 \texttt{1.csv} 自带的风格先验；最终模型用 PLS2、Ridge、ExtraTrees、RandomForest 与 LightGBM 的 OOF 等权融合，并在训练折内用保序回归纠正预测收缩后做稳健 min-max 校准。重复 $5\times3$ 折交叉验证给出宏平均 $R^2=0.227$、Spearman $=0.432$；保序校准后的 QWK 为 0.376；30\% MCAR 下 $R^2$ 相对下降 19.3\%，MAR 下下降 10.9\%。

针对问题二，本文以 \texttt{2.csv} 构建 231 个化合物的共表达网络（Top-10 近邻，8 个社区，$Q=0.586$），用边际相关、弹性网、PLS、树模型与 LightGBM SHAP 多方法融合得到纯数据 Top-3，并用 bootstrap 入选频率（均值 0.525）与方法间 Jaccard（0.313）量化稳定性。由于纯数据 Top-3 在二乙酰、DMS、T2N 上选到的是风格/物化代理而非同名直接化合物，本文另给出机制导向 Top-3（文献先验 $\lambda=0.5$）并做 $\lambda$ 敏感性分析，把“预测性关联”与“化学机制”分开报告。

针对问题三，本文在 log1p 空间用条件高斯 copula 学习 22 个风格的化合物分布，并用高斯最优传输做跨风格迁移。最终模型对 DMS/Diacetyl/4VG/T2N 采用风格感知上限 $\max(\text{阈值},\text{风格 p75})$ 做投影；三风格平均 MMD 0.038、覆盖度 0.86，Pils/Lager $\to$ Tripel 的 Tripel 分类命中率从 0.00 提升到 0.91；严格阈值迁移下四类缺陷的超严格阈值率均为 0，生成模型在 style cap 下的四类缺陷超标率均为 0（严格阈值按风格取舍另行报告）；Tripel/Kriek/Lambic 的感官画像 Pearson 分别为 0.968/0.957/0.990。

最后，本文讨论了三问结果对啤酒风味设计与质量控制的启示：化合物模型适合做配方筛选与关键化合物监测，缺陷控制仍需前置到原料与发酵过程；外部数据只能做风格级校验，不能替代逐酒感官标签。

\vspace{0.8em}
\noindent\textbf{关键词：}啤酒风味；缺失数据；多任务回归；共表达网络；稳定性选择；条件生成模型；风味指纹


\newpage

% ============================================================
% 第 1 章 问题重述
% ============================================================

\section{问题重述}

\subsection{问题背景}

啤酒的风味、口感与香气由其化学组成决定。发酵过程中，酿酒酵母的糖代谢与酯类、高级醇合成，酒花添加带来的萜类化合物，以及麦芽和微生物代谢产生的醛酮类物质，共同塑造了啤酒的果香、醇厚感、苦味与纯净度\citep{vanderhaegen2006chemistry}。同一化合物在不同浓度区间内可能呈现完全不同的感官贡献；多种化合物之间又存在显著共表达或共同代谢关系。因此，仅靠单变量分析难以刻画“化合物谱—感官画像”的整体联系。

附件给出了 250 个商用啤酒样本的化合物丰度、化合物相关矩阵、感官评分卡、感官相关矩阵以及综合参考数据。数据同时具有高维、小样本、缺失、异常值与多源异质的特点：化合物丰度部分列经过对数变换、部分列保留线性浓度，感官评分以 0--6 离散等级记录，部分表格只给出标准化后的连续值，跨表连接还存在标识列不一致的问题。上述特点决定了本文必须先完成严格的数据对齐与预处理，再针对三问分别设计模型。

\subsection{数据说明}

附件共 8 份数据，其内容与本文用途如表~\ref{tab:data-overview} 所示。

\begin{table}[htbp]
  \centering
  \caption{附件数据一览与本文用途}
  \label{tab:data-overview}
  \small
  \begin{tabularx}{\textwidth}{@{}l X X@{}}
    \toprule
    文件 & 主要内容 & 本文用途 \\
    \midrule
    \texttt{1.csv} & 250 个啤酒样本 $\times$ 化合物丰度矩阵；含少量理化参数，化合物列逐列核对（37 列 log10、33 列线性、161 列无校准按 log10）。 & 问题一至问题三的主特征表。 \\
    \texttt{2.csv} & 化合物两两之间的 Pearson 相关矩阵。 & 问题二共表达图的边权来源。 \\
    \texttt{3.xlsx} & Aroma 与 Flavour 两套感官评分卡，评分档位为 0--6。 & 问题一的目标维度定义。 \\
    \texttt{4.csv} & 感官评分的标准化/中心化连续值。 & 问题一、二的监督信号。 \\
    \texttt{5.csv} & 感官属性之间的相关矩阵。 & 感官维度结构分析与共线性诊断。 \\
    \texttt{6.csv} & 化合物/理化指标与感官属性的交叉相关矩阵。 & 问题二的外部一致性校验。 \\
    \texttt{7.csv} & 250 个样本的感官描述符，含风格标签 \texttt{tasting\_category\_fine}。 & 问题三的风格标签；问题一的辅助监督。 \\
    \texttt{8.xlsx} & 感官阈值、感官名称、GBM/SHAP、校准化学数据、价格等。 & 阈值参照、变换核验与外部校验。 \\
    \bottomrule
  \end{tabularx}
\end{table}

需要特别说明两点。第一，\texttt{4.csv} 与 \texttt{7.csv} 的 \texttt{beer\_id} 存在科学计数法存储，唯一值数量与样本数不一致；本文跨表连接统一使用 \texttt{beer} 列，不使用 \texttt{beer\_id}。第二，\texttt{3.xlsx} 的 beer 评分卡是空白评分卡而非已填写的 0--6 分数，因此问题一不存在可直接监督的等级标签；本文把 \texttt{4.csv}、\texttt{7.csv} 的标准化连续值作为连续监督信号，再通过有序校准得到 0--6 分级结果，并在结果中明确这一限制。

\subsection{问题提出}

\begin{enumerate}
  \item \textbf{问题一：数据缺失条件下的感官评分。}以 \texttt{1.csv} 的化合物数据为基础，给出 \texttt{3.xlsx} 中定义的 Aroma 与 Flavour 两套、共约 40 个感官维度的 0--6 分制分级评分结果。数据存在缺失值与异常值，且没有现成的等级标签可供直接训练。
  \item \textbf{问题二：化合物—风味分析。}利用 \texttt{2.csv} 的化合物相关数据构建化合物共表达图，结合 \texttt{1.csv} 特征，给出每个感官维度的 Top-3 关键因素化合物，并从 bootstrap 重采样、不同建模方法一致性以及化学代谢路径三个角度论证结果的稳定性与合理性。
  \item \textbf{问题三：风味与“风味指纹”生成。}以 \texttt{tasting\_category\_fine} 为风格标签，学习风格条件生成模型：给定风格（如 Tripel、Kriek、Lambic）生成合理的化合物谱与感官画像；并完成跨风格迁移（如把 Pils/Lager 的化合物谱迁移为 Tripel 风格），同时保证 DMS、双乙酰、4VG、反式-2-壬烯醛、stale hops 等缺陷化合物不恶化。
\end{enumerate}

\subsection{建模路线与逻辑递进}

三问共享同一套数据底座与符号体系，但任务性质不同：问题一是“从化合物到感官”的有监督映射，问题二是“从相关结构到关键因素”的关联分析，问题三是“从风格标签到化合物谱”的条件生成。为避免一上来堆砌复杂模型，本文对每一问都按“基线 $\rightarrow$ 改进 $\rightarrow$ 最终”组织：先给一个可复现、可解释的简单基线，再针对基线的具体缺陷逐层改进，最后用消融与稳定性实验说明每一层改进的必要性。总体路线见表~\ref{tab:roadmap} 与图~\ref{fig:roadmap}。

\begin{table}[htbp]
  \centering
  \caption{三问的递进式建模路线}
  \label{tab:roadmap}
  \small
  \begin{tabularx}{\textwidth}{@{}l X X X@{}}
    \toprule
    问题 & 基线 & 改进 & 最终 \\
    \midrule
    问题一 & 逐维均值/岭回归 & 缺失指示 + 风格先验的多输出回归 & 五模型等权融合 + 保序校准 \\
    问题二 & 边际相关排序 & 图拉普拉斯平滑的重要性 & 多方法集成 + bootstrap 稳定性选择 \\
    问题三 & 风格均值/协方差高斯采样 & 条件高斯 copula & OT + style cap 缺陷投影 \\
    \bottomrule
  \end{tabularx}
\end{table}

\figauto[1.0\textwidth]{fig_roadmap.pdf}{本文总体建模路线}{fig:roadmap}

全文安排如下：第~\ref{sec:preprocess} 节给出数据对齐、缺失处理、变换核验与符号约定；第~\ref{sec:q1} 节至第~\ref{sec:q3} 节分别回答问题一至问题三；第~\ref{sec:conclusion} 节总结结论、建议与局限；附录给出代码结构与复现步骤。

% ============================================================
% 第 2 章 数据预处理
% ============================================================

\section{数据预处理}
\label{sec:preprocess}

\subsection{符号约定}

为避免各问各自造符号，本文全文只使用一套符号，首次出现时定义，后文直接复用；角标最多两级（如 $x_{i,j}$），不引入三层及以上的上下标。主要符号见表~\ref{tab:notation}。

\begin{longtable}{@{}p{0.19\textwidth} p{0.49\textwidth} p{0.24\textwidth}@{}}
  \caption{全文符号约定}\\
  \toprule
  符号 & 含义 & 维度/取值 \\
  \midrule
  \endhead
  $N,P,S,G$ & 样本数、化合物数、感官维度数、风格类别数 & 正整数 \\
  $i,j,s,g$ & 样本、化合物、感官维度、风格索引 & $i=1,\dots,N$；其余同理 \\
  $k$ & 感官等级索引 & $k=0,1,\dots,6$ \\
  $\vx_i$ & 第 $i$ 个样本的化合物特征向量 & $\R^{P}$ \\
  $x_{i,j}$ & 第 $i$ 个样本第 $j$ 个化合物的特征值（逐列 log10 或线性，见变换表） & $\R$ \\
  $\bX$ & 样本—化合物特征矩阵 & $\R^{N\times P}$ \\
  $r_{j,j'}$ & 化合物 $j$ 与 $j'$ 的 Pearson 相关系数 & $[-1,1]$ \\
  $\bR$ & 化合物相关矩阵 & $\R^{P\times P}$ \\
  $a_{j,j'}$ & 共表达图邻接矩阵元素 & $\{0,1\}$ \\
  $\bA$ & 共表达图邻接矩阵 & $\{0,1\}^{P\times P}$ \\
  $y_{i,s}$ & 第 $i$ 个样本第 $s$ 个感官维度的评分 & $\{0,\dots,6\}$ 或连续潜变量 \\
  $\vy_i$ & 第 $i$ 个样本的感官评分向量 & $\R^{S}$ \\
  $\bY$ & 样本—感官评分矩阵 & $\R^{N\times S}$ \\
  $z_i$ & 第 $i$ 个样本的风格标签 & $\{1,\dots,G\}$ \\
  $\Dcal_g$ & 属于风格 $g$ 的样本集合 & $\Dcal_g\subseteq\{1,\dots,N\}$ \\
  $\bbeta_s$ & 第 $s$ 个感官维度的回归系数向量 & $\R^{P}$ \\
  $\bB$ & 各感官维度回归系数组成的矩阵 & $\R^{P\times S}$ \\
  $w_{j,s}$ & 化合物 $j$ 对感官维度 $s$ 的重要性 & $\R$ \\
  $\hat{y}_{i,s}$ & 模型预测的感官评分 & $\R$ \\
  $\tilde{\vx}$ & 生成的化合物谱 & $\R^{P}$ \\
  $\kappa_{s,k}$ & 感官维度 $s$ 的第 $k$ 个有序阈值 & $\R$ \\
  $\tau_j$ & 化合物 $j$ 的缺陷阈值 & $\R$ \\
  $q_{g,j}$ & 风格 $g$ 中化合物 $j$ 的 75\% 分位数 & $\R$ \\
  $\pi_{j,s}$ & 化合物 $j$ 在感官维度 $s$ 的 bootstrap 入选频率 & $[0,1]$ \\
  $\btheta$ & 模型参数 & — \\
  $\Lcal$ & 损失函数 & — \\
  \bottomrule
\end{longtable}
\label{tab:notation}

\subsection{数据概况与表间对齐}

附件各表的样本标识并不完全一致。\texttt{1.csv}、\texttt{4.csv}、\texttt{7.csv} 均含 \texttt{beer} 列，取值为 1--250 的唯一编号；\texttt{4.csv} 与 \texttt{7.csv} 的 \texttt{beer\_id} 存在大量科学计数法存储，无法作为可靠主键。因此本文以 \texttt{beer} 为唯一连接键，先建立样本主表，再逐表左连接，连接后核对每列的非空数量与取值范围。

样本主表共 250 行、231 个化合物特征。自然缺失 73 个单元（0.12\%），主要集中在 \texttt{amyl\_acetate}（27）、\texttt{aroma\_hops\_sum}（13）、\texttt{methylpyrazine}（12）、\texttt{2,3-butanedione}（9）与 \texttt{3-methylbutanal}（7）。感官目标共 50 维，列标准化后的标准差为 0.37--0.86，目标间相关中位数为 0.016，但存在若干强相关对（最高 $r=0.98$），说明不能把 50 个维度当作完全独立的任务。风格标签共 22 类，最大类 Blond 31 个样本、Tripel 28 个，最小类 Faro 3 个；小样本风格在问题三中采用收缩估计与 bootstrap 报告不确定性。

化合物丰度的变换并不是整表统一的。将 \texttt{8.xlsx} 的 Calibrated chemical dataset 与 \texttt{1.csv} 的 70 个公共化合物列逐列比较：37 列满足 $\text{calibrated}=10^{x_{i,j}}$（log10），33 列满足 $\text{calibrated}\approx x_{i,j}$（线性浓度）；其余 161 列没有校准浓度，按负值与值域判定为 log10。本文据此建立逐列变换表，只在解释浓度时按列还原，建模时统一用分位数变换，避免把 $10^{x}$ 误套到线性列上。完整的 231 列变换表随支撑材料提供，正文只报告 37/33/161 的统计与典型列（表~\ref{tab:transform-summary}）。

\begin{table}[htbp]
  \centering
  \caption{\texttt{1.csv} 逐列变换核对结果}
  \label{tab:transform-summary}
  \small
  \begin{tabular}{@{}l c l l@{}}
    \toprule
    类型 & 列数 & 判定依据 & 典型列 \\
    \midrule
    校准 log10 & 37 & $\text{calibrated}=10^{x}$ & \texttt{ethyl\_acetate}、\texttt{DMS}、\texttt{pH} \\
    校准线性 & 33 & $\text{calibrated}\approx x$ & \texttt{linalool}、\texttt{geraniol}、\texttt{2,3-butanedione} \\
    推断 log10 & 161 & 负值或值域 $<12$ & \texttt{X4VG}、萜类、酯类 \\
    \bottomrule
  \end{tabular}
\end{table}

感官数据方面，\texttt{3.xlsx} 只提供评分卡结构与 0--6 档位，不提供已填写的分数；\texttt{4.csv} 与 \texttt{7.csv} 给出的是标准化/中心化后的连续值，含负值与缺失。两者不能直接当作 0--6 标签使用。本文的处理是：以连续值作为监督信号训练回归或有序模型，再通过稳健区间映射与分位数校准得到 0--6 分级；所有涉及“真实等级”的指标仅在模拟等级或专家复核数据上报告，并在文中注明。

\texttt{4.csv}（训练面板）与 \texttt{7.csv}（Ratebeer 描述符）并非同一套感官体系。对可对应的概念逐对计算相关：酸感 0.824、酒精感 0.653、苦味 0.653、甜感 0.573，说明两者在主要口感维度上外部一致；但 DMS 仅 0.119、反式-2-壬烯醛 0.194、双乙酰 $-0.009$，说明痕量缺陷维度不能靠 \texttt{7.csv} 校验。本文因此把 \texttt{7.csv} 用作主要口感维度的外部稳健性检查，而不把它当作缺陷标签。

进一步用 Kaggle 的 22 个风格原型做独立校验：它与 \texttt{7.csv} 的风格均值在酸感（$r=0.890$）、麦芽（0.798）、酒花（0.819）、苦味（0.748）、酒精（0.688）、甜感（0.697）上方向一致；果味（0.147）与酒体（$-0.267$）则不可靠。与 \texttt{4.csv} 的相关整体更弱（酒精 0.262、麦芽 0.429、酒花 0.449），说明训练面板与评论词频并不是同一套感官体系。外部数据只用于风格级排序校验与先验，不能当作个体级 0--6 标签；这一限制在第~\ref{sec:q3} 节生成评价中同样适用。

\subsection{缺失值与异常值处理}

\texttt{1.csv} 与 \texttt{7.csv} 均存在缺失。对化合物矩阵，本文用缺失掩码 $\bM\in\{0,1\}^{N\times P}$ 记录观测位置，$m_{i,j}=1$ 表示第 $i$ 个样本第 $j$ 个化合物被观测到。缺失处理遵循两条原则：其一，任何插补、标准化与阈值选择都只在训练折内拟合，验证/测试折只做变换，防止信息泄漏；其二，插补方法作为敏感性分析的一个因素，而不是只保留一种“最优”结果。

基线插补采用训练集逐列中位数；Q1 另加缺失指示特征，Q3 对 \texttt{X\_cal} 的 3 个负值截断为 0（\texttt{betaglucan.mg.L.} 与 \texttt{unfermentables}）。异常值未做统一删除：化合物侧的极端值多为真实风味差异，删除会损失信息；本文在 Q1 的分位数变换、Q3 的 log1p 与分位数校准中降低其影响，并在结果中保留极端样本。缺失数据的统计处理框架参照 Little 与 Rubin 的讨论\citep{little2019statistical}。

\subsection{变换、标准化与特征构造}

化合物特征以 log10 为主，但 33 个校准列保留线性浓度；逐列变换后，不同化合物量纲差异仍然很大，本文对特征做训练集内标准化：
\begin{equation}
  \tilde{x}_{i,j}=\frac{x_{i,j}-\mu_j}{\sigma_j},
  \label{eq:standardize}
\end{equation}
其中 $\mu_j$、$\sigma_j$ 为训练集第 $j$ 列的均值与标准差。对组成型数据，本文另做中心对数比（CLR）变换作为敏感性分析，检查结论是否依赖单一变换。

特征构造遵循“可解释优先”的原则：保留原始化合物列作为主特征；按 \texttt{8.xlsx} 的 Sensory names 与化合物类别，把化合物归入酒花萜类、酵母酯类、高级醇类、醛酮类、麦芽来源与其他六组，构造组内均值与总和特征；对缺陷化合物单独建组，便于问题三的约束与问题二的合理性解释。所有派生特征在论文中注明来源列与计算方式，不使用无法追溯的临时特征。

\subsection{外部数据与补充信息}

题目给出的 Kaggle Top Beer Information 数据集\citep{kaggle2026topbeer} 已实际下载并核对：主表 5,558 行 $\times$ 21 列，来源为 BeerAdvocate 评论词频，许可为 CC BY 4.0。它与附件 250 款酒没有共同的酒名、酒厂或条码字段，无法逐酒连接；本文只在 22 个 \texttt{tasting\_category\_fine} 风格层级使用它：将 112 个 Kaggle 风格中的 96 个映射到附件 22 个风格，覆盖 4,760/5,558 行（85.6\%）。具体用途有三：其一，为 Q1 提供风格级感官先验与外部稳健性校验；其二，为 Q3 提供 22 个风格的感官原型与迁移方向；其三，用风格级 ABV 做交叉检查（$r=0.855$），而 IBU 的相关仅 0.449，故不把 IBU 强行用于个体样本。外部数据只做风格级辅助，不替代 \texttt{1.csv}/\texttt{2.csv} 的化合物建模。

\subsection{预处理流程与复现}

预处理流程如算法~\ref{alg:preprocess} 所示。该流程只依赖附件原始数据与本文代码，不依赖任何手工修改的中间表；随机种子、缺失率与折数统一记录在配置文件与附录中。

\begin{algorithm}[htbp]
  \small
  \caption{多源数据对齐与预处理}
  \label{alg:preprocess}
  \Input{附件 \texttt{1.csv}--\texttt{8.xlsx}}
  \Output{对齐特征矩阵 $\bX$、评分矩阵 $\bY$、风格标签 $\{z_i\}$、缺失掩码 $\bM$}
  读取各表，统一以 \texttt{beer} 为样本主键，核对行数与唯一性\;
  用 \texttt{8.xlsx} 的校准化学数据逐列核验 \texttt{1.csv} 的变换（37 列 log10、33 列线性）\;
  建立缺失掩码 $\bM$，标记 \texttt{1.csv} 与 \texttt{7.csv} 的缺失位置\;
  按样本分层划分训练/验证/测试集，固定随机种子\;
  \For{每个训练折}{
    在训练折内拟合中位数插补、均值与标准差，并变换其余各折\;
  }
  解析 \texttt{7.csv} 的 \texttt{tasting\_category\_fine}，建立风格集合 $\Dcal_g$\;
  输出对齐后的数据与一份可复查的预处理日志\;
\end{algorithm}

\figauto{fig_preprocess.pdf}{数据预处理与表间对齐流程}{fig:preprocess}

预处理流程如图~\ref{fig:preprocess} 所示：先对齐样本主键，再核验变换、建立掩码、划分数据与拟合变换参数，最后输出可复查的日志。该顺序保证任何监督信息都不会从验证/测试折泄漏到训练折。

需要强调的是，本文不把 \texttt{4.csv} 与 \texttt{7.csv} 的标准化值反标准化后直接当作 0--6 分数：标准化是在不同样本集上分别完成的，反变换会引入不可识别的一致性问题。正确做法是保留连续监督信号，另建从连续潜变量到 0--6 等级的有序校准层，详见第~\ref{sec:q1} 节。

% ============================================================
% 第 3 章 问题一：数据缺失条件下的感官评分
% ============================================================

\section{问题一：数据缺失条件下的感官评分}
\label{sec:q1}

\subsection{问题分析}

问题一要求给出 \texttt{3.xlsx} 中 Aroma 与 Flavour 两套、共约 40 个感官维度的 0--6 分评分。具体地，\texttt{3.xlsx} 的 Aroma 与 Flavour 各 15 维（共 30），\texttt{4.csv} 另含 20 个 taste/palate/overall 维度；本文对全部 50 维建模，并在结果中把 Aroma/Flavour 30 维作为核心、其余 20 维作为辅助输出。该任务有三个特殊之处。

第一，\textbf{没有可直接使用的等级标签}。\texttt{3.xlsx} 的 beer 评分卡是空白评分卡，\texttt{4.csv} 与 \texttt{7.csv} 提供的是标准化/中心化后的连续值，不是原始 0--6 分。若把标准化值直接四舍五入，会因样本集不同、中心化基准不同而产生不可比的分级。因此本文把任务拆成两层：先学习连续感官潜变量，再用有序阈值把它校准为 0--6 等级。

第二，\textbf{化合物侧存在缺失与高维共线}。\texttt{1.csv} 的化合物维度远大于样本量，且同源化合物高度相关；直接逐维回归会过拟合，缺失值处理不当还会把“缺失模式”误学成“风味信号”。因此模型必须显式使用缺失掩码，并在维度间共享信息。

第三，\textbf{感官维度之间存在相关结构}。\texttt{5.csv} 与 \texttt{6.csv} 表明不少感官维度共同变化；把 50 个维度完全独立建模会损失这部分结构。最终模型因此采用多输出/低秩共享结构，让不同感官维度共享化合物因子，并用集成融合补充非线性。

\subsection{评价指标}

设测试集共有 $n$ 个样本、$S$ 个感官维度。连续预测的评价指标为
\begin{align}
  \mathrm{MAE} &= \frac{1}{nS}\sum_{i=1}^{n}\sum_{s=1}^{S}\abs{\hat{y}_{i,s}-y_{i,s}},
  \label{eq:mae}\\
  \mathrm{RMSE} &= \sqrt{\frac{1}{nS}\sum_{i=1}^{n}\sum_{s=1}^{S}\left(\hat{y}_{i,s}-y_{i,s}\right)^2}.
  \label{eq:rmse}
\end{align}
等级预测采用二次加权 Kappa（QWK）与 Spearman 秩相关：
\begin{equation}
  \mathrm{QWK}=1-\frac{\sum_{k=0}^{6}\sum_{k'=0}^{6}W_{k,k'}O_{k,k'}}
  {\sum_{k=0}^{6}\sum_{k'=0}^{6}W_{k,k'}E_{k,k'}},
  \label{eq:qwk}
\end{equation}
其中 $O_{k,k'}$ 为真实等级 $k$ 被预测为 $k'$ 的观测频数，$E_{k,k'}$ 为对应期望频数，$W_{k,k'}=(k-k')^2/36$ 为二次权重\citep{cohen1968weighted}。Spearman 系数按
\begin{equation}
  \rho_s=1-\frac{6\sum_{i=1}^{n}\delta_i^2}{n(n^2-1)}
  \label{eq:spearman}
\end{equation}
计算，$\delta_i$ 为第 $i$ 个样本在两个排序中的名次差\citep{spearman1904proof}。

由于真实 0--6 标签不可得，本文在测试集上同时报告两类评价：其一是对连续监督信号的 MAE、RMSE 与 Spearman；其二是对“模拟等级”或专家复核等级的 QWK。前者衡量潜变量拟合，后者衡量分级校准；两者不可混为一谈。这也是本文把分级结果定位为“校准后的相对等级”而不是“绝对人工评分”的原因。

还需明确“参考分级”的含义：本文的参考分级是把连续监督信号按训练折内 $q_{0.01}/q_{0.99}$ 做稳健 min-max 后取整得到的离散结果，它是连续值的单调重标度，不含人工评分的绝对档位信息。因此后文报告的 QWK 衡量的是“同一连续信号在两种离散化下的一致性”，性质上接近离散化后的秩相关，不等于与人工评分逐档吻合。

\subsection{基线模型：逐维均值与线性回归}

最直接的基线是不使用化合物信息，只用训练集均值预测每个感官维度：
\begin{equation}
  \hat{y}_{i,s}^{(0)}=\bar{y}_s,\qquad
  \bar{y}_s=\frac{1}{N}\sum_{i=1}^{N}y_{i,s}.
  \label{eq:q1-mean}
\end{equation}
它给出“没有任何化合物信息”的下界。第二个基线是逐维岭回归，设 $\bbeta_s\in\R^{P}$ 为第 $s$ 个维度的回归系数，$c_s$ 为截距，$\bB=[\bbeta_1,\dots,\bbeta_S]$：
\begin{equation}
  \min_{\bB,\{c_s\}}
  \sum_{i=1}^{N}\sum_{s=1}^{S}\left(y_{i,s}-\vx_i^\top\bbeta_s-c_s\right)^2
  +\lambda\norm{\bB}_F^2.
  \label{eq:q1-ridge}
\end{equation}
岭回归通过共享特征矩阵与统一正则化降低小样本方差；PLS2 则把 $\bB$ 约束为低秩，用少量潜变量同时解释多个感官维度，是“多任务共享”的最简形式。二者都不使用缺失指示与风格信息，作为后续改进的对照\citep{hastie2009elements,tibshirani1996regression,wold2001pls}。

\subsection{改进模型：缺失感知的多输出回归与风格先验}

改进模型在岭回归/PLS2 的基础上加入两项数据事实。第一，化合物侧存在缺失：本文在训练折内做中位数插补，并把“某列是否缺失”作为 0/1 指示特征拼入 $\vx_i$，让模型区分“浓度低”与“未检测”。第二，\texttt{1.csv} 自带风格标签 $g(i)$，而同一化合物谱在不同风格下的感官期望并不相同；本文把风格先验拼入特征：
\begin{equation}
  \tilde{\vx}_i=\left[\vx_i;\ \bs_i\right],\qquad
  \bs_i\in\{0,1\}^{G}\ \text{或}\ \bs_i=\bar{\vy}_{g(i)},
  \label{eq:q1-style}
\end{equation}
其中 $\bar{\vy}_{g}$ 为训练折内风格 $g$ 的感官均值（target encoding），用留一或平滑避免泄漏。风格先验不是“用标签作弊”：它是 \texttt{1.csv} 的已有列，且在真实场景中风格在品评前已知；本文同时报告 chemical-only 与 chemical+style 的消融，量化它带来的提升。

\subsection{最终模型：集成融合与单调校准}

单一模型对共线化合物与非线性关系都不够稳健。最终模型先用五类基学习器得到 OOF 预测：PLS2（低秩多任务）、岭回归（线性）、ExtraTrees 与 RandomForest（非线性、多输出）、LightGBM（逐目标 boosting）。记第 $\nu$ 个基学习器对样本 $i$、维度 $s$ 的 OOF 预测为 $\hat{y}_{i,s}^{(\nu)}$，$\Mcal$ 为基学习器集合。本文比较两种融合：
\begin{equation}
  \hat{y}_{i,s}^{\mathrm{avg}}=\frac{1}{\abs{\Mcal}}\sum_{\nu\in\Mcal}\hat{y}_{i,s}^{(\nu)},
  \label{eq:q1-ensemble}
\end{equation}
\begin{equation}
  \hat{y}_{i,s}^{\mathrm{stack}}=\sum_{\nu\in\Mcal}\omega_{\nu,s}\,\hat{y}_{i,s}^{(\nu)}+b_s,
  \label{eq:q1-stack}
\end{equation}
其中 $\omega_{\nu,s}$ 在 OOF 预测上拟合。15 折结果显示等权融合的 $R^2=0.227$ 高于 Ridge 堆叠的 0.216，且折间方差更小，故最终采用式~\eqref{eq:q1-ensemble}；堆叠作为对照保留。融合的价值在于：线性模型与树模型在不同感官维度上各有胜负，等权平均让模型在不估计额外权重的情况下获得稳健性。

堆叠方案存在一处需要如实说明的乐观偏置：元学习器在训练折样本的折外预测上拟合，而这些折外预测来自训练时包含当前测试折样本的模型，故 0.216 这一数值偏高。即便带此偏置，堆叠仍未超过等权融合的 0.227，这进一步支持选择等权融合作为最终方案。

连续预测存在向均值收缩的问题，直接做 min-max 会把 0--6 分级挤到中间几档。本文在训练折内拟合单调校准 $g_s$（isotonic regression，保序回归），把 OOF 预测映射到真实连续值，再做稳健 min-max：
\begin{equation}
  \hat{k}_{i,s}=\operatorname{round}\left(
  6\,\frac{g_s(\hat{y}_{i,s})-q_{0.01,s}}{q_{0.99,s}-q_{0.01,s}}
  \right),\qquad \hat{k}_{i,s}\in\{0,\dots,6\},
  \label{eq:q1-grade}
\end{equation}
其中 $q_{0.01,s},q_{0.99,s}$ 为训练折内 $g_s(\hat{y}_{\cdot,s})$ 的 1\%/99\% 分位数，越界值截断。保序回归只要求单调，不假设线性，能纠正收缩而不改变样本排序；min-max 只负责把校准后的连续值锚定到 0--6。quantile 与 Gaussian latent 两种校准作为敏感性分析，用于检查分级结论是否依赖校准假设。

\figauto{fig_q1_framework.pdf}{问题一的分级评分框架：基线、改进与最终模型}{fig:q1framework}

图~\ref{fig:q1framework} 概括了三种路线的输入、输出与衔接关系：基线只用均值/线性回归；改进模型加入缺失指示与风格先验；最终模型用集成融合与保序校准。三者在同一数据划分上比较，避免“换了数据又换了模型”的不可比问题。

\subsection{实验设置}

本文按样本分层划分训练、验证与测试集，所有随机过程固定种子；缺失插补、标准化、保序校准与超参数选择只在训练/验证集内完成。为检验缺失条件下的稳健性，在测试集上按不同比例随机遮蔽化合物，并额外模拟分位数驱动的非随机缺失（MAR），比较三种路线的退化曲线。主要超参数包括正则化强度 $\lambda$、PLS 潜变量数、树模型叶子数/树数、LightGBM 学习率、集成权重与校准分位数；其搜索范围与最终取值记录在附录表~\ref{tab:appendix-config} 中。本文使用 5 折 $\times$ 3 次重复（种子 42），MCAR 5/10/20/30\%，MAR 由 \texttt{ethanol..v.v.} 的分位数驱动。所有超参数均在外层训练折内用内层交叉验证选择、不跨折固定：岭回归 $\alpha\in[10^{0},10^{6}]$（13 点对数网格、内层 5 折）、PLS 成分数 $2$--$20$（内层 5 折）、ExtraTrees 与 RandomForest 叶节点 $\in\{1,2,4\}$（内层 3 折），树模型最终 500 棵；LightGBM 固定 400 棵、学习率 0.03、叶子 15。堆叠元学习器为 RidgeCV（$\alpha\in[10^{-3},10^{2}]$、内层 3 折）。各折实际取值记录在 \texttt{data/q1/model\_summary.json} 的 \texttt{hyperparameters} 字段。

对照设置包括：（1）逐维均值基线；（2）逐维岭回归；（3）式~\eqref{eq:q1-ridge} 的多输出岭回归/PLS2；（4）式~\eqref{eq:q1-ensemble} 的等权融合；（5）去掉风格先验、去掉保序校准、或去掉树模型的消融版本。通过这一组对照，可以分辨“多任务共享”“风格先验”和“保序校准”各自的贡献，而不是只报一个最终分数。

\subsection{结果与分析}

表~\ref{tab:q1-results} 给出三种路线在重复交叉验证上的连续拟合与分级一致性指标。最终模型相对基线 Ridge 的 $R^2$ 提升 0.065（0.162$\to$0.227），Spearman 提升 0.059（0.373$\to$0.432）。

\begin{table}[htbp]
  \centering
  \caption{问题一三种路线的交叉验证结果}
  \label{tab:q1-results}
  \small
  \begin{tabular}{@{}l c c c c c@{}}
    \toprule
    模型 & $R^2$ & MAE & RMSE & Spearman $\rho_s$ & QWK \\
    \midrule
    基线：逐维岭回归 & 0.1617 & 0.3622 & 0.4946 & 0.3728 & — \\
    改进：岭回归 + 风格先验 & 0.1855 & 0.3542 & 0.4855 & 0.3927 & — \\
    最终：五模型等权融合 + 保序校准 & 0.2270 & 0.3415 & 0.4695 & 0.4318 & 0.3756 \\
    \bottomrule
  \end{tabular}
\end{table}

表中 R²/RMSE/MAE/Spearman 为 5 折 $\times$ 3 次重复的折间均值；QWK 只对最终分级校准报告，基线与改进未做分级校准。最终模型的折间标准差为 R² 0.021、Spearman 0.024、RMSE 0.018、MAE 0.011；基线与改进的标准差见支撑材料。风格先验在岭回归上稳定提升 R² $+0.024$（0.1617$\to$0.1855），在融合模型上提升 $+0.012$（0.2270$\to$0.2388）；由于融合模型已通过化学特征间接吸收风格，最终交付仍以 chemical-only 融合为主，style 版本作为消融报告。

从指标结构看，连续指标与分级指标回答的是不同问题：MAE/RMSE 衡量连续潜变量是否贴合监督信号，QWK 衡量分级是否与等级顺序一致。最终模型的 14/50 个目标 $R^2>0.3$，其中 acidity（0.690）、alcohol（0.678）、lactic（0.648）、acetic（0.635）、bitternes（0.604）最强；barnyard（$-0.039$）、stale\_hops（$-0.032$）、clove（$-0.005$）最弱。复杂/微生物/阈值型维度基本不可由 231 个化学特征预测，这与化学合理性文献一致，也解释了为什么 Q2 的纯数据 Top-3 会在这些维度上选到代理。

\figauto{q1_target_r2_heatmap.pdf}{逐感官维度的交叉验证 $R^2$}{fig:q1target}

在缺失稳健性方面，自然缺失（73/57,750 $=0.126\%$）下最佳模型的 $R^2$ 为 0.2264；MCAR 5\%/10\%/20\%/30\% 时 $R^2$ 相对下降 3.46\%/6.33\%/11.46\%/19.33\%，RMSE 相对上升 0.70\%/1.15\%/2.28\%/3.75\%；MAR（24.77\%）的 $R^2$ 下降 10.89\%，与同量级 MCAR 接近。30\% MCAR 内退化可控，说明缺失指示与折内中位数插补有效；但没有证据表明模型能完全消除非随机缺失偏差，这一限制在讨论中保留。

\figauto{q1_missingness.pdf}{缺失率—性能曲线（MCAR 与 MAR）}{fig:q1missing}

\subsection{讨论}

问题一的主要困难不是模型复杂度，而是标签不可直接观测。本文用“连续监督 + 有序校准”绕开了空白评分卡的限制，但这也意味着最终 0--6 分数的绝对水平依赖校准锚点，不能声称与人工评分逐档完全一致。更稳妥的表述是：在给定校准方案下，模型能保持样本间的相对顺序与等级单调性；若要报告绝对等级精度，需要原始评分或专家复核。

其次，低秩结构假定不同感官维度由少量共享因子驱动，这与 \texttt{5.csv}、\texttt{6.csv} 的相关结构一致，但若某些维度由独立机制主导（例如特定缺陷阈值），低秩约束可能欠拟合。解决办法是允许部分维度使用独立残差项，或在低秩模型与逐维模型之间做模型平均；这部分作为后续改进方向。

最后，缺失机制未知。本文的掩码损失在随机缺失下是无偏的，在非随机缺失下只能减轻偏差，不能消除。若后续能获得缺失原因的元数据（例如某化合物未检测而非浓度为零），应把缺失原因作为额外变量纳入模型。

% ============================================================
% 第 4 章 问题二：化合物—风味关联分析
% ============================================================

\section{问题二：化合物—风味关联分析}
\label{sec:q2}

\subsection{问题分析}

问题二要求为每个感官维度给出 Top-3 关键因素化合物，并论证结果的稳定性与化学合理性。它的核心不是“训练一个预测精度最高的模型”，而是“在大量高度共线的化合物中，找出对某个感官维度既有统计支持、又有化学解释的少数变量”。

附件提供了三类可用信息：\texttt{2.csv} 的化合物—化合物相关矩阵、\texttt{1.csv} 的样本级化合物丰度、\texttt{6.csv} 的化合物—感官交叉相关矩阵。三者分别刻画“化合物之间是否共同变化”“样本间差异是否与感官相关”“已有研究/数据中是否观察到对应关系”。一个稳健的 Top-3 应该在三类信息上都不出现明显矛盾，而不是只在一张相关矩阵上排到前列。

因此本文把问题二组织为三步：先由 \texttt{2.csv} 与 \texttt{1.csv} 构建共表达图，刻画化合物之间的模块结构；再在图上做重要性排序，得到每个感官维度的候选化合物；最后用多方法集成、bootstrap 重采样与化学类别核对筛选 Top-3，并明确报告稳定性指标。

\subsection{共表达图构建}

设化合物 $j$ 与 $j'$ 的相关系数为 $r_{j,j'}$（来自 \texttt{2.csv}），邻接矩阵 $\bA$ 的元素定义为
\begin{equation}
  a_{j,j'}=
  \begin{cases}
    1, & \abs{r_{j,j'}}\ge \eta,\\[2pt]
    0, & \text{否则},
  \end{cases}
  \qquad a_{j,j}=0,
  \label{eq:q2-graph}
\end{equation}
其中阈值 $\eta$ 不是固定常数，而是在验证集上按“图连通性—模块可解释性—下游稳定性”三者折中选取。若简单阈值导致孤立点过多，则改用每个节点保留 Top-$K$ 强相关的 $K$ 近邻图，并在附录中报告两种图的差异。

图的拉普拉斯矩阵定义为
\begin{equation}
  \bL=\bD-\bA,\qquad
  d_j=\sum_{j'=1}^{P}a_{j,j'},
  \label{eq:q2-laplacian}
\end{equation}
其中 $\bD=\operatorname{diag}(d_1,\dots,d_P)$ 为度矩阵。对 $\bL$ 做特征分解可以得到化合物的低频模块结构；本文用模块度与社区划分检查图是否把酒花萜类、酵母酯类、高级醇类等已知类别分开。共表达网络的构建与模块分析参照 WGCNA 的做法\citep{langfelder2008wgcna}。

\figauto[1.0\textwidth]{q2_network.pdf}{化合物共表达网络与社区结构示例}{fig:q2network}

图~\ref{fig:q2network} 给出共表达网络的示意结构：节点为化合物，边表示强相关，颜色表示社区。网络既是后续图平滑重要性的输入，也是化学合理性论证的可视化依据。主方案（每个节点保留 Top-10 强相关邻居）得到 231 个节点、1,683 条边、8 个社区，模块度 $Q=0.586$；社区大小与主要化学类别见表~\ref{tab:q2-communities}。与硬阈值 $|r|\ge0.5$（48 个孤立点）和 Top-5 近邻（865 条边）相比，Top-10 在保留全部节点的同时维持了清晰的模块结构，故作为主方案。

\begin{table}[htbp]
  \centering
  \caption{共表达网络的 8 个社区及其化学解释}
  \label{tab:q2-communities}
  \small
  \begin{tabularx}{\textwidth}{@{}l c l l X@{}}
    \toprule
    社区 & 节点数 & 主要来源 & 主要类别 & 代表成员 \\
    \midrule
    C1 & 40 & 酵母/酒花 & 酯/萜/醇/醛 & E-nerolidol、2,3-dihydrofarnesol、n-dodecanol \\
    C2 & 34 & 酒花 & 萜类/酯 & nerol oxide、$\alpha$-calacorene、epicubenol \\
    C3 & 32 & 酒花/酵母/野生菌 & 酯/挥发性酚 & ethyl phenylacetate、$\beta$-ionone、neryl formate \\
    C4 & 29 & 酵母 & 酯/醇/硫/理化 & ethanol、total esters、isopentyl acetate \\
    C5 & 28 & 麦芽/酵母 & 酯/醛/呋喃吡嗪 & diacetyl、$\beta$-farnesene、2-pentylfuran \\
    C6 & 24 & 酵母/酒花 & 酯（16/24） & 2-methylbutyl octanoate、isopentyl hexanoate \\
    C7 & 22 & 酵母/酒花 & 醇/酯/萜 & 3-furanmethanol、3-methyl-2-hexanol、o-cresol \\
    C8 & 22 & 酒花/香料 & 萜类（11/22） & isoborneol、terpinolene、terpinen-4-ol \\
    \bottomrule
  \end{tabularx}
\end{table}

\subsection{基线方法：边际相关排序}

最直接的 Top-3 是逐化合物计算其与目标感官维度的边际相关，再取绝对值最大的三个：
\begin{equation}
  w_{j,s}^{(0)}=\abs{\Corr\left(x_{\cdot,j},y_{\cdot,s}\right)},
  \qquad
  \Tcal_s^{(0)}=\operatorname*{Top3}_{j}\ w_{j,s}^{(0)},
  \label{eq:q2-baseline}
\end{equation}
其中 $x_{\cdot,j}$ 为第 $j$ 个化合物的样本向量，$y_{\cdot,s}$ 为第 $s$ 个感官维度的评分向量。该基线完全可复现、解释直接，但存在三个问题：其一，共线化合物会“轮流”占据前列，Top-3 不稳定；其二，边际相关无法区分直接关联与通过第三变量产生的间接关联；其三，缺失值与异常值会显著影响单个相关系数。下面的改进与最终模型分别针对这三点。

\subsection{改进方法：图平滑重要性}

共表达图提供了“哪些化合物应该一起被考虑”的先验。本文把基线重要性向量 $\bw_s^{(0)}$ 在图上做拉普拉斯平滑：
\begin{equation}
  \tilde{\bw}_s=\left(\bI+\alpha\bL\right)^{-1}\bw_s^{(0)},
  \label{eq:q2-smooth}
\end{equation}
其中 $\alpha>0$ 控制平滑强度。式~\eqref{eq:q2-smooth} 的直观含义是：一个化合物的重要性不仅取决于它自己的边际相关，还取决于与它强连接的邻居；孤立的高相关点会被适度压低，成模块出现的高相关化合物会被整体保留。与直接对相关系数做阈值相比，这种做法不切断弱相关边，而是让重要性沿图扩散，因而对阈值选择不那么敏感。

平滑强度 $\alpha$ 由验证集上的 Top-3 稳定性选择，而不是固定为 1。若 $\alpha\to 0$，模型退化为基线；若 $\alpha$ 过大，所有化合物的重要性会被拉平，Top-3 失去区分度。这一“从基线出发、由数据选择改进幅度”的设计，使改进模型始终包含基线作为特例。

\subsection{最终方法：多方法集成与稳定性选择}

单一建模方法对共线结构敏感，因此最终模型并列保留五种互补方法，并在排序层面做秩融合：

\begin{enumerate}
  \item \textbf{A 边际 Spearman}：逐化合物计算与目标感官维度的 Spearman 相关，在每个目标内做 BH-FDR 校正，按 $|r|$ 排序；
  \item \textbf{B 弹性网}：对每个感官维度做弹性网回归，内层 4 折交叉验证在 $\alpha$ 与 $l_1$ 比例网格上选参，用系数绝对值作为重要性\citep{zou2005regularization}；
  \item \textbf{C 极端随机树}：逐目标拟合 300 棵树（深度 8、特征子采样 $\sqrt{P}$、叶节点 3），用不纯度重要性\citep{breiman2001random}；
  \item \textbf{D LightGBM SHAP}：逐目标做 5 折折外 LightGBM（150 棵树），用折外平均 $|\mathrm{SHAP}|$ 衡量化合物对单个感官维度的贡献\citep{friedman2001greedy,chen2016xgboost,lundberg2017unified}；
  \item \textbf{E 网络平滑}：把 A 的重要性沿共表达图做一步邻居平滑（$\lambda=0.30$）。
\end{enumerate}

记方法集合为 $\Mcal=\{A,B,C,D,E\}$，第 $\nu$ 种方法给出的重要性为 $\tilde{w}_{j,s}^{(\nu)}$，各自在每个目标内排秩得 $R^{(\nu)}_{j,s}$。融合采用 Borda 秩平均：
\begin{equation}
  \bar{R}_{j,s}=\frac{1}{|\Mcal|}\sum_{\nu\in\Mcal}R^{(\nu)}_{j,s},
  \qquad
  w_{j,s}^{*}=1-\frac{\bar{R}_{j,s}-1}{P-1},
  \label{eq:q2-ensemble}
\end{equation}
其中 $w_{j,s}^{*}\in[0,1]$，最终 Top-3 取 $\bar{R}_{j,s}$ 最小的三个化合物。C 与 D 只产生无符号重要性，融合前用 A 的边际相关符号为其赋号，得到带方向的共识分数。等权秩平均不设可调权重，避免在小样本验证集上调权导致二次过拟合。

需要说明两处口径差异。其一，方法 D 在主分析用 5 折折外 SHAP，在 bootstrap 中为控制耗时改用全样本 SHAP，因此表中 \texttt{rank\_D} 反映折外贡献、\texttt{freq\_D} 反映样本内贡献，二者不可逐项对照，比较应以融合结果为准。其二，E 只改变同一目标内重要性的分布形状，不引入新的数据信息，故仅作为家族感知的审计视角。

稳定性用 bootstrap 入选频率与 Jaccard 系数量化。对第 $s$ 个感官维度，重复 $B$ 次有放回重采样，记第 $b$ 次得到的 Top-3 集合为 $\Tcal_s^{(b)}$，则化合物 $j$ 的入选频率为
\begin{equation}
  \pi_{j,s}=\frac{1}{B}\sum_{b=1}^{B}\mathbf{1}\left\{j\in\Tcal_s^{(b)}\right\}.
  \label{eq:q2-frequency}
\end{equation}
两种方法 $\nu$ 与 $\nu'$ 之间的一致性为
\begin{equation}
  J\left(\Tcal_s^{(\nu)},\Tcal_s^{(\nu')}\right)=
  \frac{\abs{\Tcal_s^{(\nu)}\cap\Tcal_s^{(\nu')}}}
  {\abs{\Tcal_s^{(\nu)}\cup\Tcal_s^{(\nu')}}}.
  \label{eq:q2-jaccard}
\end{equation}
最终 Top-3 取 $\bar{R}_{j,s}$ 最小的三个化合物，并同步报告其 bootstrap 入选频率 $\pi_{j,s}$；入选频率低于 0.5 的条目在表中标注“低稳定”，但不因此替换或补足，以保证 Top-3 完全由点估计决定、稳定性作为独立信息报告\citep{efron1979bootstrap,jaccard1901etude}。

\begin{algorithm}[htbp]
  \small
  \caption{基于 bootstrap 的 Top-3 稳定性选择}
  \label{alg:stability}
  \Input{训练集特征 $\bX$、感官矩阵 $\bY$、共表达图 $\bA$、方法集合 $\Mcal$、重采样次数 $B$}
  \Output{每个感官维度的 Top-3 集合与入选频率 $\pi_{j,s}$}
  \For{$b=1$ \KwTo $B$}{
    从训练集中有放回抽取样本，得到第 $b$ 个重采样集\;
    \For{$\nu\in\Mcal$}{
      用方法 $\nu$ 拟合模型，计算并平滑重要性 $\tilde{w}_{j,s}^{(\nu)}$\;
    }
    按式~\eqref{eq:q2-ensemble} 集成并记录每个感官维度的 Top-3 集合 $\Tcal_s^{(b)}$\;
  }
  按式~\eqref{eq:q2-frequency} 计算 $\pi_{j,s}$，与最终 Top-3 一并报告（不按频率二次筛选）\;
  计算方法间 Jaccard 系数，输出稳定性报告\;
\end{algorithm}

\subsection{实验设置}

本文把样本按 \texttt{beer} 分层划分训练/验证/测试集，所有方法在同一划分上比较；缺失值在训练折内插补，测试折只做变换。主要超参数包括图阈值 $\eta$（或近邻数 $K$）、平滑强度 $\alpha$、弹性网的混合比例与正则化强度、PLS 潜变量数、树模型深度与学习率；搜索范围记录在附录表~\ref{tab:appendix-config}。稳定性实验重采样次数 $B$ 与入选频率门槛也在附录中给出。主方案 $K=10$；图平滑 $\alpha=0.3$；Elastic Net 用内层 CV 选 $\alpha$ 与 $l_1$ 比例；PLS 用内层 CV 选成分数；树模型 300 棵（bootstrap 60 棵）；LightGBM 150 棵、5 折 OOF SHAP；bootstrap $B=500$；入选频率仅用于报告稳定性、不参与筛选；机制先验 $\lambda=0.5$。

评价分为四层：（1）预测层：各方法对感官维度的交叉验证误差；（2）排序层：Top-3 的 bootstrap 入选频率；（3）一致层：方法间 Jaccard 系数；（4）化学层：Top-3 是否落在合理的化合物类别、方向是否与已知代谢路径一致。只有四层都不出现明显矛盾，才把该组 Top-3 作为最终结果。

\subsection{结果与分析}

表~\ref{tab:q2-top3} 对 8 个代表性维度同时给出纯数据 Top-3 与机制导向 Top-3。纯数据 Top-3 由 \texttt{2.csv}/\texttt{1.csv} 的统计关联决定，回答“哪些化合物最能预测该感官”；机制导向 Top-3 在融合分数上加一个文献先验（$\lambda=0.5$，直接化合物权重 1.0），回答“哪些化合物在化学上最应被关注”。两者不是替代关系：纯数据结果反映预测性关联，机制导向结果用于化学审计；把两者混为一谈会把相关当因果。

\begin{table}[htbp]
  \centering
  \caption{纯数据 Top-3 与机制导向 Top-3 的对照（8 个代表性维度）}
  \label{tab:q2-top3}
  \footnotesize
  \setlength{\tabcolsep}{3pt}
  \begin{tabularx}{\textwidth}{@{}l X X X@{}}
    \toprule
    感官维度 & 纯数据 Top-3 & 机制导向 Top-3（$\lambda=0.5$） & 说明 \\
    \midrule
    diacetyl & Pentyl octanoate; Ethyl decanoate; $\beta$-phellandrene & Diacetyl; Pentyl octanoate; Ethyl decanoate & 纯数据选酯/萜烯代理；机制导向把二乙酰放回第 1 \\
    dms & Protein; Color; total sulfur & Total sulfur; SO$_2$; DMS & 纯数据选蛋白/色度；机制导向把 DMS 放回第 3 \\
    t2n & Glycerol; $\alpha$-phellandrene; $\beta$-phellandrene & E-2-nonenal; dimethoxy-E-citral; Z-dihydro-apofarnesal & 纯数据无醛类；机制导向选醛类 \\
    X4vg & 4-vinylguaiacol; neryl formate; prenyl isobutyrate & 4-vinylguaiacol; 4-ethyl-guaiacol; 4-ethyl-phenol & 两视图一致：挥发性酚 \\
    A\_esters\_fruity & $\beta$-phellandrene; epicubenol; humulene epoxide II & Neryl formate; Ethyl decanoate; 2-phenylethyl butyrate & 纯数据选萜烯；机制导向选酯类 \\
    A\_hops\_citrus & Total hop aroma; cis-dihydro-$\beta$-terpineol; citronellol & Citronellol; Geraniol; Humulene epoxide II & 两视图都落在萜类/酒花 \\
    A\_malt\_all & Color; Caloric value; pH & 2-acetylfuran; 5-methylfurfural; Furfural & 纯数据选物化代理；机制导向选麦芽美拉德产物 \\
    bitternes & Ethyl octanoate; Iso-alpha acids; Ethyl decanoate & Iso-alpha acids; Ethyl octanoate; Ethyl decanoate & 两视图一致：异 $\alpha$-酸 + 酯类 \\
    \bottomrule
  \end{tabularx}
\end{table}

对二乙酰、DMS、T2N 三个缺陷维度，纯数据 Top-3 没有选到同名直接化合物：二乙酰的直接化合物 fused rank=120.8，DMS=86.5，T2N=91.4；机制导向在 $\lambda=0.5$ 时把它们放回第 1/3/1 位。这个差异本身是重要发现：在本数据中，直接缺陷化合物的边际信号弱于风格/物化代理，可能来自阈值效应、测量噪声与风格混杂；因此不能仅凭相关排序断言“二乙酰是二乙酰感官的最强预测因子”。相反，X4vg 与 bitternes 的两视图一致，说明这两条关联既有统计支持又有化学对应。

$\lambda$ 的敏感性见表~\ref{tab:q2-domain-sensitivity}：$\lambda=0$ 时机制导向退化为纯数据，Top-3 重合率 1.0；$\lambda=0.5$ 时重合率 0.387，9 个直接/特定化合物进入 Top-3；$\lambda$ 增大不改变与 \texttt{6.csv} 的符号一致率（1.0），但 Top-3 与 \texttt{6.csv} 的重合率从 0.46 降到 0.31。$\lambda=0.5$ 不是最优，而是“强先验”情景；本文不以它替代纯数据结果，只把它作为机制审计。还需注意量纲：融合分数 $w_{j,s}^{*}\in[0,1]$，而机制先验加分在 $\lambda=0.5$、直接化合物权重 1.0 时达 0.5，与数据分数同量级。因此 $\lambda=0.5$ 下同名直接化合物进入第 1 位在很大程度上是加分幅度的算术后果，只能说明“在该先验强度下必然如此”，不能作为数据支持该化合物的证据。

\begin{table}[htbp]
  \centering
  \caption{机制先验强度 $\lambda$ 的敏感性}
  \label{tab:q2-domain-sensitivity}
  \footnotesize
  \setlength{\tabcolsep}{3pt}
  \begin{tabular}{@{}c c c c c@{}}
    \toprule
    $\lambda$ & 与纯数据 Top-3 重合率 & 直接/特定化合物数 & 6.csv 符号一致率 & 6.csv Top-3 重合率 \\
    \midrule
    0.0 & 1.000 & 5 & 1.0 & 0.460 \\
    0.1 & 0.780 & 5 & 1.0 & 0.393 \\
    0.2 & 0.633 & 5 & 1.0 & 0.387 \\
    0.5 & 0.387 & 9 & 1.0 & 0.313 \\
    \bottomrule
  \end{tabular}
\end{table}

结果分析区分三类：第一类是多方法一致、bootstrap 高频、化学路径清楚，例如 bitternes 的 iso-alpha acids、X4vg 的 4-vinylguaiacol、A\_hops\_citrus 的 citronellol/geraniol；第二类是统计上稳定但化学解释弱，例如 dms 的 protein/color、A\_malt\_all 的 color/caloric value/pH，这类标注为数据驱动关联，并说明可能是风格/工艺混杂；第三类是只在机制先验下进入 Top-3 的直接化合物（diacetyl、DMS、T2N），作为机制审计线索，不声称它们是最强预测因子。这样的分层报告比单纯列三个名字更符合题目“稳定性与化学合理性论证”的要求。

\subsection{稳定性与化学合理性论证}

\paragraph{稳定性。}500 次 bootstrap 下，纯数据 Top-3 的平均入选频率为 0.525，中位数 0.47；150 个条目中 29 个达到 0.8、67 个达到 0.5。没有感官维度的三个条目全部达到 0.8，6 个维度三个条目都达到 0.5，24 个维度至少两个条目达到 0.5，22 个维度至少一个条目达到 0.8；方法间 Jaccard 均值为 0.313。这说明“单个化合物稳定”比“整组 Top-3 稳定”更常见，Top-3 应带着频率报告，而不是只给三个名字。

\figauto{q2_top3_stability.pdf}{纯数据 Top-3 的 bootstrap 入选频率}{fig:q2stability}

\figauto[1.0\textwidth]{q2_jaccard.pdf}{五种方法之间的 Top-3 Jaccard 系数}{fig:q2jaccard}

图构造方案的扰动也支持这一点：Top-5 与 Top-10 的社区划分 ARI 为 0.558，阈值 $|r|\ge0.5$ 会产生 48 个孤立点，因此本文选 Top-10 作为主方案，并在图~\ref{fig:q2network} 中标注社区。

\paragraph{与 \texttt{6.csv} 的关系。}\texttt{6.csv} 的 44 个化合物与 \texttt{2.csv}/\texttt{1.csv} 同源；方法 A 的边际相关排序与 \texttt{6.csv} 的平均秩相关接近 1.0。这验证了表间对齐，但不能当作独立复现；多变量方法（B/C/D/E）与 \texttt{6.csv} 的偏离是预期的，因为 \texttt{6.csv} 是边际相关，未控制共线。方法 E 在 $\lambda=0.3$ 时把秩相关从 1.0 降到 0.982、Top-3 重合率从 1.0 降到 0.82，因此只作为 family-aware 审计，不作为“更好”的证据。

\paragraph{化学合理性。}网络社区与已知化学家族对齐：C2/C8 富集酒花萜类，C4/C6 富集酵母酯类/醇类，C5 富集麦芽醛/呋喃/吡嗪（表~\ref{tab:q2-communities}）。Top-3 中，bitternes 的 iso-alpha acids、X4vg 的 4-vinylguaiacol、A\_hops\_citrus 的 citronellol/geraniol 与文献一致；diacetyl/DMS/T2N 的直接化合物在纯数据中排名靠后，但在机制先验下进入 Top-3，说明它们的作用可能被风格/阈值效应掩盖。类别富集在 Top-3 上有 3 个显著、Top-10 上有 9 个显著（$q<0.05$），数量不多，说明“类别富集”只能作为辅助证据，不能替代逐化合物核对。

这里还有两点必须澄清，否则会高估“两视图一致”的证据强度。第一，\texttt{6.csv} 的 44 行中“Iso-alpha acids”与 \texttt{1.csv} 的 \texttt{bitterness} 列同源，而异 $\alpha$-酸正是苦味物质的化学来源，因此 \texttt{bitternes} 维度上的“两视图一致”部分属于定义性一致，不构成独立证据；其余维度不受此影响。第二，\texttt{1.csv} 末尾的 \texttt{esters\_sum}、\texttt{aroma\_hops\_sum}、\texttt{acids\_sum}、\texttt{sulfur\_sum} 由组分相加得到，本文把它们与单个化合物同样视为节点，这会人为提高其枢纽程度、并使酯类与酒花类目标更容易选到汇总列，解释 Top-3 时应与具体组分对照阅读。

\paragraph{与问题一的衔接。}问题二得到的 Top-3 可以作为问题一模型的可解释性检查：若某个化合物在问题二中被判为某感官维度的关键因素，却在问题一的多任务模型中没有非零或稳定系数，说明两者至少有一个需要复查。本文不把这种不一致藏起来，而是把它作为模型诊断信息报告。

\subsection{讨论}

问题二的结论受三类因素限制。第一，相关不等于因果：共表达图上的边可能来自共同上游调控、测量批次或样品制备，而非直接化学作用；本文用多方法与 \texttt{6.csv} 交叉校验降低这种风险，但不能证明因果。第二，化合物注释与类别划分依赖 \texttt{8.xlsx} 与文献，若注释有误，化学合理性判断会跟着错；因此本文把类别核对过程与来源写进附录，便于复查。第三，Top-3 是“给定数据与给定感官维度下的紧凑摘要”，不是“只有这三个化合物重要”；排名第 4、5 位的化合物可能在特定风格或特定缺陷场景中更关键，实际使用时不应把 Top-3 当作硬阈值。

% ============================================================
% 第 5 章 问题三：风味指纹与风格条件生成
% ============================================================

\section{问题三：风味指纹与风格条件生成}
\label{sec:q3}

\subsection{问题分析}

问题三包含三个相互关联的子任务：其一，学习以 \texttt{tasting\_category\_fine} 为条件的化合物谱生成模型，使给定风格（如 Tripel、Kriek、Lambic）能生成合理的化合物组合；其二，由生成的化合物谱给出对应的感官画像，即“风味指纹”；其三，完成跨风格迁移，例如把 Pils/Lager 的化合物谱迁移为 Tripel 风格，同时保证 DMS、双乙酰、4VG、反式-2-壬烯醛与 stale hops 等缺陷化合物不恶化。

与问题一、二相比，问题三的困难在于“没有唯一的正确答案”。同一风格的不同啤酒在化合物谱上差异很大，生成模型不能只记住训练样本；同时，风格标签在 \texttt{7.csv} 中分布不均，部分风格样本很少，直接对每个风格单独建模会过拟合。与 GAN/CVAE 类神经生成方法\citep{goodfellow2014generative,mirza2014conditional} 相比，条件高斯 copula 在 $N=250$、$G=22$ 的小样本下参数更少、协方差可收缩，更适合本题；本文因此采用“条件 copula 学风格分布、高斯最优传输做跨风格迁移、style cap 做缺陷约束”的路线，并把评价拆成分布保真、风格保真、多样性、缺陷控制与感官一致五个维度，避免只看生成样本“像不像”。

\subsection{风格标签与风味指纹定义}

以 \texttt{7.csv} 的 \texttt{tasting\_category\_fine} 作为风格标签 $z_i$，通过 \texttt{beer} 列与 \texttt{1.csv} 连接，得到化合物谱—风格配对数据。风格集合记为 $\{1,\dots,G\}$，属于风格 $g$ 的样本集合为 $\Dcal_g$。由于各风格样本量差异较大，本文在训练时采用风格分层划分与协方差收缩，并在结果中单独报告小样本风格的表现。

本文把风格 $g$ 的“风味指纹”定义为其化合物谱均值与感官画像均值的组合：
\begin{equation}
  \mathcal{F}_g=\left(\bar{\vx}_g,\bar{\vy}_g\right),\qquad
  \bar{\vx}_g=\frac{1}{\abs{\Dcal_g}}\sum_{i\in\Dcal_g}\vx_i,\quad
  \bar{\vy}_g=\frac{1}{\abs{\Dcal_g}}\sum_{i\in\Dcal_g}\vy_i.
  \label{eq:q3-fingerprint}
\end{equation}
该定义不是要替代完整的生成分布，而是给生成结果提供一个可比较的“风格中心”：生成样本既要接近 $\bar{\vx}_g$，又要保留风格内部的合理变异，不能退化为均值复制。

\subsection{基线模型：风格均值与全局协方差}

最简单的风格条件生成器是“风格均值 + 高斯噪声”：对每个风格 $g$，估计均值 $\bmu_g$，并用全局协方差 $\bSigma_{\mathrm{pool}}$ 注入随机性：
\begin{equation}
  \tilde{\vx}^{(g)}=\bmu_g+\bSigma_{\mathrm{pool}}^{1/2}\bz,\qquad
  \bz\sim\Ncal(\bzero,\bI).
  \label{eq:q3-mean-noise}
\end{equation}
第二个基线是“全局协方差 + 风格均值平移”：保留全局相关结构，只把均值移到目标风格。它们可复现、训练快，能给出风格中心与基本变异，但无法刻画风格内协方差的差异，也容易在长尾/零膨胀化合物上生成不合理的极端值；下面的条件高斯 copula 针对这一点改进。

\subsection{改进模型：条件高斯 copula}

化合物浓度具有强零膨胀与长尾，直接用高斯模型会生成负值或极端值。本文先把每个化合物 $j$ 的经验分布 $F_j$ 映射到标准正态分数
\begin{equation}
  z_{i,j}=\Phi^{-1}\left(F_j(x_{i,j})\right),
  \label{eq:q3-normal-score}
\end{equation}
其中 $F_j$ 由训练集的经验分布（或分位数变换）估计，$\Phi$ 为标准正态累积分布函数。对风格 $g$，估计均值 $\bmu_g$ 与协方差 $\bSigma_g$；当 $\abs{\Dcal_g}$ 较小时用收缩估计
\begin{equation}
  \bSigma_g^{(\lambda)}=(1-\lambda)\hat{\bSigma}_g+\lambda\bSigma_{\mathrm{pool}},
  \label{eq:q3-shrinkage}
\end{equation}
其中 $\bSigma_{\mathrm{pool}}$ 为所有风格合并后的协方差，$\lambda\in[0,1]$ 由交叉验证或 Ledoit--Wolf 规则选择。生成时
\begin{equation}
  \bz\sim\Ncal(\bmu_g,\bSigma_g^{(\lambda)}),\qquad
  \tilde{x}_j=F_j^{-1}\left(\Phi(z_j)\right),
  \label{eq:q3-copula}
\end{equation}
并对 $\tilde{x}_j$ 做非负截断。式~\eqref{eq:q3-copula} 是一个条件高斯 copula：它保留每个化合物的边际分布与风格内的相关结构，同时把风格信息压缩为均值与协方差。与风格均值/全局协方差基线相比，收缩协方差在小样本风格上更稳定；与条件变分自编码器（CVAE）相比，它没有神经网络的高方差，更适合 $N=250$、$G=22$ 的数据规模。本文把 CVAE 作为后续扩展，而不是当前主模型\citep{kingma2014auto,sohn2015learning}。

\subsection{最终模型：缺陷约束生成与最优传输迁移}

\paragraph{缺陷约束。}记缺陷化合物集合为 $\Jcal$（至少包含 DMS、双乙酰、4VG、反式-2-壬烯醛；stale hops 用代理变量），其阈值为 $\tau_j$。定义风格感知上限
\begin{equation}
  \bar{c}_{g,j}=\max\left(\tau_j,\ q_{g,j}\right),
  \label{eq:q3-cap}
\end{equation}
其中 $q_{g,j}$ 为风格 $g$ 中化合物 $j$ 的 75\% 分位数。生成后对缺陷列做投影
\begin{equation}
  \tilde{x}_j\leftarrow\min\left(\tilde{x}_j,\bar{c}_{g,j}\right),\qquad j\in\Jcal.
  \label{eq:q3-projection}
\end{equation}
这样既不把缺陷一律压到零（避免破坏风格基线），也不允许生成超出风格正常波动的极端值；同时报告严格阈值 $\tau_j$ 下的超阈值率，区分“风格允许”与“绝对超标”。

其中 4VG 的风格上限取固定安全值 $0.10$~mg/L 而非 $\max(\tau_j,q_{g,j})$，因为小麦风格 4VG 的 $p_{75}$ 显著高于该安全值；其余三个缺陷按式~\eqref{eq:q3-cap} 取风格上限。stale hops 不是单一化合物，本文用反式-2-壬烯醛、DMS、\texttt{ethyl 2-methylbutyrate}、\texttt{ethyl isovalerate} 与 4VG 五列标准化均值构造代理，并在迁移中对每一列施加“不高于源样本”的约束。各缺陷的阈值与依据见表~\ref{tab:q3-thresholds}。

\begin{table}[htbp]
  \centering
  \caption{缺陷化合物的严格阈值与取值依据}
  \label{tab:q3-thresholds}
  \footnotesize
  \setlength{\tabcolsep}{3pt}
  \begin{tabular}{@{}l l l l l@{}}
    \toprule
    缺陷 & 对应列 & 严格阈值 $\tau_j$ & 单位 & 依据 \\
    \midrule
    DMS & \texttt{DMS} & 17.25 & \textmu g/L & 附件与文献综合 \\
    Diacetyl & \texttt{X2.3.butanedione} & 0.09 & mg/L & 附件与文献综合 \\
    4VG & \texttt{X4VG} & 0.10 & mg/L & 非小麦安全值（均值 0.206） \\
    trans-2-nonenal & \texttt{X.E..2.nonenal} & 0.10 & \textmu g/L & 外部文献 \\
    stale hops & 五列复合代理 & — & — & 见正文说明 \\
    \bottomrule
  \end{tabular}
\end{table}

\paragraph{跨风格迁移。}设源风格为 $g$、目标风格为 $g'$。在正态分数空间，用高斯最优传输的闭式映射
\begin{equation}
  T_{g\to g'}(\bz)=\bmu_{g'}+\bSigma_{g'}^{1/2}\bSigma_g^{-1/2}\left(\bz-\bmu_g\right),
  \label{eq:q3-ot}
\end{equation}
将源样本映射到目标风格，再逆变换回浓度空间并做式~\eqref{eq:q3-projection} 的缺陷投影。式~\eqref{eq:q3-ot} 是均值和协方差匹配意义下的最优传输映射；它比“只加风格均值差”多了协方差校正，能同时改变化合物的相关结构。对 Pils/Lager $\to$ Tripel，本文比较三种迁移：raw（只做 OT，不加缺陷约束）、strict（投影到附件严格阈值）、style（投影到式~\eqref{eq:q3-cap} 的风格上限），并报告三者的缺陷超标率与感官画像。

\paragraph{目标边际校准。}仿射 OT 只匹配均值与协方差，不保证各化合物的边际形状（零膨胀的痕量化合物尤甚）。因此迁移后对每个化合物列再做一次秩映射，把生成样本的分位数对齐到目标风格的经验分布：
\begin{equation}
  \tilde{x}_{i,j}\leftarrow Q_{g',j}\!\left(\frac{\operatorname{rank}(\tilde{x}_{i,j})-1/2}{n}\right),
  \label{eq:q3-rank}
\end{equation}
其中 $Q_{g',j}$ 为目标风格 $g'$ 化合物 $j$ 的经验分位数函数。该步之后每一列的边际与目标风格严格一致，因此后文的 MMD 与覆盖度衡量的是\textbf{化合物之间的相关结构}，而非整体分布形状；引用 $0.038$ 这一数值时必须带上此口径。

\begin{algorithm}[htbp]
  \small
  \caption{缺陷约束的最优传输风格迁移}
  \label{alg:transfer}
  \Input{源样本 $\vx_i$、源风格 $g$、目标风格 $g'$、边际分布 $F_j$、阈值 $\tau_j$、风格分位数 $q_{g',j}$}
  \Output{迁移后的化合物谱 $\tilde{\vx}_{\mathrm{tr}}$}
  计算源/目标风格的正态分数均值与协方差 $\bmu_g,\bSigma_g,\bmu_{g'},\bSigma_{g'}$\;
  对源样本做正态分数变换 $\bz_i=\Phi^{-1}(F(\vx_i))$\;
  用式~\eqref{eq:q3-ot} 映射到目标风格，再逆变换回浓度空间\;
  \For{每个缺陷化合物 $j\in\Jcal$}{
    投影 $\tilde{x}_j\leftarrow\min(\tilde{x}_j,\max(\tau_j,q_{g',j}))$\;
  }
  输出 $\tilde{\vx}_{\mathrm{tr}}$\;
\end{algorithm}

\subsection{实验设置与评价指标}

本文按风格分层划分训练/验证/测试集，保证每个风格在测试集中至少有一个样本；小样本风格采用收缩估计，并在结果中单独标注其不确定性。生成模型的对照包括：风格均值 + 高斯噪声（mean\_noise）、全局协方差 + 风格均值平移（global\_cov）、条件高斯 copula、copula + 缺陷约束，以及 OT 迁移。所有模型使用相同的训练/测试划分与随机种子。Tripel 28、Kriek 10、Lambic 11 个样本；每风格每模型生成 300 条；输入为 70 个校准化合物，做 log1p 与逐列标准化；协方差用 Ledoit--Wolf 收缩；OT 后做逐特征目标分位数校准；缺陷投影用 $\max(\tau_j,q_{g,j})$；感官代理用 11 成分 PLS。

评价从五个维度展开：

\begin{enumerate}
  \item \textbf{风格保真}：用真实数据训练风格分类器，在生成样本上的分类准确率；同时计算生成样本到目标风格中心的马氏距离
  \begin{equation}
    D_g(\tilde{\vx})=\sqrt{\left(\tilde{\vx}-\bmu_g\right)^\top
    \bSigma_g^{-1}\left(\tilde{\vx}-\bmu_g\right)}.
    \label{eq:q3-mahalanobis}
  \end{equation}
  \item \textbf{多样性}：生成样本两两之间的平均距离，以及各化合物维度的方差；防止模型只输出风格均值。另有覆盖度与精确度两个派生指标：以真实样本留一最近邻距离的中位数为阈值，覆盖度定义为“至少被一个生成样本落入该阈值的真实样本比例”；以真实留一最近邻距离的 95\% 分位为阈值，精确度定义为“到真实集最近邻距离不超过该阈值的生成样本比例”。
  \item \textbf{缺陷控制}：缺陷化合物超阈值比例
  \begin{equation}
    R_{\mathrm{def}}=\frac{1}{\abs{\Jcal}}
    \sum_{j\in\Jcal}\mathbf{1}\left\{\tilde{x}_j>\tau_j\right\}.
    \label{eq:q3-defect-rate}
  \end{equation}
  \item \textbf{感官一致性}：用问题一最终模型预测生成谱的感官画像 $\hat{\vy}(\tilde{\vx})$，与目标风格画像 $\bar{\vy}_g$ 比较
  \begin{equation}
    \mathrm{MAE}_{\mathrm{sen}}=\frac{1}{S}\sum_{s=1}^{S}
    \abs{\hat{y}_s(\tilde{\vx})-\bar{y}_{g,s}}.
    \label{eq:q3-sensory}
  \end{equation}
\end{enumerate}

此外，用 RBF 核的最大均值差异（MMD）衡量生成分布与真实风格分布的整体差异：
\begin{equation}
  \mathrm{MMD}^2=
  \frac{1}{n^2}\sum_{i,i'}k(\vx_i,\vx_{i'})
  +\frac{1}{m^2}\sum_{j,j'}k(\tilde{\vx}_j,\tilde{\vx}_{j'})
  -\frac{2}{nm}\sum_{i,j}k(\vx_i,\tilde{\vx}_j),
  \label{eq:q3-mmd}
\end{equation}
其中 $k(\cdot,\cdot)$ 为 RBF 核。五个维度必须同时报告：只提高风格分类准确率可能意味着模式崩塌，只提高多样性可能生成不合理化合物，只压低缺陷比例可能把风味一起压没。本文用“分布—保真—多样—缺陷—感官”五象限表呈现结果，避免单一指标误导。

\subsection{结果与分析}

表~\ref{tab:q3-results} 给出三风格平均的五维评价。OT+约束在分布保真（MMD 0.038）、覆盖度（0.863）与 top-1 命中率（0.981）上均最好；条件 copula 的命中率 0.974 接近，但 MMD（0.111）与感官 RMSE（0.132）更差，说明它更接近风格中心、却牺牲了分布与感官一致性；均值噪声的命中率 0.851、覆盖度 0.300，存在模式崩塌。缺陷列显示，OT+约束把 T2N/4VG 的超严格阈值率降到 0，但 DMS/Diacetyl 仍分别有 0.59/0.81 的超严格阈值率，这是因为 Tripel/Kriek/Lambic 的风格 p75 本身高于附件严格阈值；相对 style cap 的超标率为 0（四缺陷均 0.000），详见迁移段落。因此“风格允许”与“绝对安全”必须分开报告；表中的超严格阈值率按 DMS/Diacetyl/T2N 三项报告。

\figauto{q3_mmd_classifier.pdf}{生成模型的分布距离与风格分类命中率}{fig:q3mmd}

\figauto[1.0\textwidth]{q3_transfer_defects.pdf}{Pils/Lager $\to$ Tripel 迁移的缺陷控制}{fig:q3transfer}

迁移结果方面，源风格 Pils/Lager 的 Tripel 分类命中率为 0.00；OT 迁移后 raw/strict/style 三种变体的 Tripel 命中率均为 0.91（随机森林），感官画像与 Tripel 的 Pearson 相关分别为 0.974/0.856/0.972，说明风格方向迁移成功。raw 迁移把 DMS/Diacetyl/T2N 从 15.8 \textmu g/L、0.117 mg/L、0.004 \textmu g/L 提高到 41.8、0.418、0.115，11/11 个样本均“恶化”；strict 投影把三者压到 13.7、0.047、0.004，超严格阈值率为 0，且相对源样本的恶化比例为 0；style 投影保留 Tripel 的风格基线（39.8、0.398、0.099），相对 Tripel p75 的超标率为 0，但相对 Pils 的升高是风格本身的特征，不应误判为缺陷恶化。关键化合物均值从 Pils 的乙醇 5.19\%、甘油 1.64 g/L、乙酸乙酯 10.8 mg/L、乙酸异戊酯 1.15 mg/L、异戊醇 27.9 mg/L，迁移到 Tripel 的 8.79\%、3.11 g/L、40.3 mg/L、4.62 mg/L、54.6 mg/L；linalool 均值从 0.062 升到 1.14 mg/L，高于 Tripel 中位数 0.372 mg/L 但低于其均值 2.97 mg/L，说明 OT 在长尾化合物上可能过冲，后续应加分位数约束。整体上，“绝对阈值”与“风格基线”两种口径给出的结论不同：前者更保守，后者更符合风格真实性；本文同时报告两者，不把“风格允许”误写成“绝对安全”。

\begin{table}[htbp]
  \centering
  \caption{问题三生成模型的五维评价（三风格平均）}
  \label{tab:q3-results}
  \footnotesize
  \setlength{\tabcolsep}{3pt}
  \begin{tabularx}{\textwidth}{@{}l X c c c c@{}}
    \toprule
    评价维度 & 指标 & 均值噪声 & 全局协方差 & 条件 copula & OT+约束 \\
    \midrule
    分布保真 & $\mathrm{MMD}^2$（越低越好） & 0.063 & 0.066 & 0.111 & 0.038 \\
    风格保真 & 随机森林 top-1 命中率 & 0.851 & 0.667 & 0.974 & 0.981 \\
    多样性 & 覆盖度（越高越好） & 0.300 & 0.233 & 0.430 & 0.863 \\
    缺陷控制 & 超严格阈值率 & 0.61/0.67/0.40 & 0.57/0.60/0.43 & 0.56/0.73/0.07 & 0.59/0.81/0.00 \\
    感官一致 & 与目标画像的 $\mathrm{RMSE}$ & 0.090 & 0.095 & 0.132 & 0.090 \\
    \bottomrule
  \end{tabularx}
\end{table}

\figauto[1.0\textwidth]{q3_radar.pdf}{各风格风味指纹的感官画像（雷达图示例）}{fig:q3radar}

图~\ref{fig:q3radar} 用雷达图对比各风格的感官画像。雷达图的每个轴对应一个感官维度，取值来自问题一模型的预测或真实样本均值；它用于检查生成样本是否在多个感官轴上同时接近目标风格，而不是只在单一轴上“像”。

案例部分重点展示三类结果：其一，给定 Tripel 生成的代表性化合物谱与感官画像，检查其是否呈现预期的果香/醇厚特征；其二，Kriek 与 Lambic 的对比，检查模型是否区分了风格间的关键差异，而不是只学到“都是酸啤”；其三，Pils/Lager 到 Tripel 的迁移，检查缺陷化合物是否仍在阈值内、感官画像是否向 Tripel 移动。案例结果为：Tripel/Kriek/Lambic 的生成画像与真实风格中位数的 Pearson 分别为 0.968/0.957/0.990，RMSE 0.051/0.147/0.071；Kriek 的 RMSE 较大，反映其真实感官分布更分散。Pils $\to$ Tripel 的 strict 迁移在保持 Tripel 命中率 0.91 的同时把四类缺陷超严格阈值率压到 0，style 迁移保留 0.972 的感官相关但保留风格基线缺陷；两者分别对应“安全优先”和“风格优先”两种使用场景。

\subsection{讨论}

问题三的评价没有唯一真值，本文用五维指标与案例结合，但仍有三点局限。第一，风格分类器本身是在真实数据上训练的，若真实数据存在标签噪声，生成样本的“保真”分数会继承这种噪声；因此本文同时报告马氏距离与感官一致性，避免只依赖分类器。第二，缺陷阈值来自 \texttt{8.xlsx} 与题目给定口径，不同风格、不同市场的阈值可能不同；本文把阈值作为可配置参数，并在敏感性分析中改变阈值，检查结论是否稳健。第三，OT 映射假设正态分数空间中的风格差异可由均值/协方差线性刻画；若风格流形高度非线性，线性 OT 可能只迁移了部分特征。后续可改用核 OT 或条件流模型，但必须保留缺陷约束与可解释性检查。

另有两处口径必须在结论前交代。其一，\texttt{X.E..2.nonenal} 在附件中标注为 mg/L，但其数值范围（Pils/Lager 中位数约 0.002、Tripel 约 0.011）与文献常用的 \textmu g/L 阈值量级一致，本文按 \textmu g/L 解释；若按标注单位计算，全部样本都远低于检测限，分类将失去意义。其二，stale hops 代理中的 \texttt{ethyl 2-methylbutyrate} 与 \texttt{ethyl isovalerate} 本身属于果香酯类，把它们纳入“不恶化”约束会同时抑制 Tripel 应有的酯类上升，因此 strict 迁移对这两列是\textbf{保守偏置}，不能读作对 stale hops 的独立实测。

另一个需要如实说明的问题是数据规模。部分风格样本很少，生成模型在该风格上的多样性指标可能被高估或低估；本文不把少数风格的结果外推到所有商业啤酒，而是把它们作为“数据条件下的可行方案”，并建议后续补充同风格样本后再复评。

% ============================================================
% 第 6 章 结论与建议
% ============================================================

\section{结论与建议}
\label{sec:conclusion}

\subsection{主要结论}

本文围绕啤酒化合物—感官关系，完成了从数据对齐、分级评分、关键化合物识别到风格条件生成的三段式建模。主要结论可概括如下。

\textbf{问题一：}在缺少原始 0--6 等级标签的条件下，本文先用逐维均值/岭回归建立下界，再用缺失指示、分位数变换与风格先验改进多输出回归，最后用 PLS2、Ridge、ExtraTrees、RandomForest 与 LightGBM 的 OOF 等权融合与保序校准输出 0--6 分级。重复 $5\times3$ 折交叉验证显示，等权融合模型的宏平均 $R^2=0.227$、Spearman $=0.432$，保序校准后的 QWK 为 0.376；在 30\% MCAR 下 $R^2$ 相对下降 19.3\%，MAR 下下降 10.9\%，说明缺失指示与折内插补有效；保序校准把 0 级预测比例从 raw min-max 的 7.6\% 提到 11.3\%，明显缓解了分级塌缩，但仍低估参考分级的尾部（参考分级中 0 级占 23.3\%），这一限制在讨论中保留。

\textbf{问题二：}纯数据 Top-3 的平均 bootstrap 入选频率为 0.525、方法间 Jaccard 为 0.313；机制导向 Top-3 在 $\lambda=0.5$ 时把二乙酰、DMS、T2N 的直接化合物放回第 1/3/1 位。两者差异说明：在本数据中，缺陷直接化合物的边际信号弱于风格/物化代理，相关排序不能直接当作因果；共表达网络的 8 个社区与酒花萜类、酵母酯类、麦芽醛/呋喃/吡嗪对齐（$Q=0.586$），提供了独立于单变量相关的化学结构证据。

\textbf{问题三：}条件高斯 copula 与高斯最优传输在 log1p 空间建模；OT+约束在三风格上平均 MMD 为 0.038、覆盖度为 0.863，优于均值噪声（0.063/0.300）、全局协方差（0.066/0.233）与 copula（0.111/0.430）。Pils/Lager $\to$ Tripel 的 Tripel 分类命中率从 0.00 提升到 0.91；迁移后严格阈值下 DMS/Diacetyl/4VG/T2N 的超严格阈值率均为 0，style-p75 迁移在保留 Tripel 风格基线（DMS 39.8 \textmu g/L、Diacetyl 0.398 mg/L）的同时把 above-style-cap 率控制在 0；生成模型的 style cap 超标率四缺陷均为 0，严格阈值超标率按风格取舍另行报告。感官代理模型的嵌套 CV 平均 $R^2=0.220$，缺陷维度的预测较弱，这是本文明确报告的局限。

\subsection{对酿造实践的建议}

第一，把“化合物—感官”模型作为配方与工艺调整的辅助工具，而不是替代人工品评。模型可以快速筛出可能影响目标风味的化合物，但最终判断仍需结合感官小组与生产条件。

第二，对关键化合物建立分风格的监控窗口。问题二得到的 Top-3 可以作为优先监测项，但应结合问题一的维度权重与问题三的缺陷约束，避免只盯少数化合物而忽略整体风味平衡。

第三，缺陷控制应前置到原料与发酵过程。DMS、双乙酰、4VG、反式-2-壬烯醛与 stale hops 等指标与原料新鲜度、煮沸强度、发酵温度与储运条件密切相关；生成模型的阈值约束只能用于配方筛选，不能替代工艺控制。

第四，数据积累应优先补齐“缺失原因”和“原始 0--6 评分”。本文的校准层是对缺失标签的补救，不是对真实等级的替代；若能获得原始评分卡与检测限信息，问题一的绝对等级精度与问题三的感官评价都会显著改善。

\subsection{局限与展望}

本文的局限主要有四方面。其一，问题一缺少真实 0--6 标签，分级结果的绝对水平依赖校准假设；本文通过敏感性分析与有序约束降低风险，但不能消除。其二，问题二的关联分析基于观测数据，相关不等于因果；共表达图上的模块可能混合了共同调控、批次效应与测量误差。其三，问题三的生成质量评价缺少统一真值，本文的五维指标与案例只能提供相对比较，不能证明生成样本在真实酿造中可行。其四，外部数据只能整合到 22 个风格层级：Kaggle 数据与附件没有共同的酒名、酒厂或条码，无法逐酒连接；本文用它做风格级原型与外部效度校验，不能替代个体级标签\citep{kaggle2026topbeer}。

后续工作可以从三个方向展开：一是引入原始感官评分与检测限数据，把问题一的分级校准从“分位数锚定”升级为“有序标签监督”；二是用干预实验或代谢网络先验替代纯相关图，提升问题二因果解释的可信度；三是把问题三的生成模型与工艺参数连接，形成“工艺—化合物—感官”的闭环优化，而不是只在化合物层面做条件生成。


\bibliographystyle{plainnat}
\bibliography{refs}

\appendix
% ============================================================
% 附录：代码与复现
% ============================================================

\section{代码与复现}

\subsection{目录结构}

本文的论文源码、代码、图片与原始数据分开存放，目录职责如表~\ref{tab:appendix-layout} 所示。原始数据只读，任何派生数据与中间结果都不写回 \texttt{附件/}。

\begin{table}[htbp]
  \centering
  \caption{项目目录与职责}
  \label{tab:appendix-layout}
  \small
  \begin{tabularx}{\textwidth}{@{}l X@{}}
    \toprule
    路径 & 职责 \\
    \midrule
    \texttt{latex/} & 论文源码：入口、导言区、参考文献与分节文件。 \\
    \texttt{latex/main.tex} & 论文入口，按顺序读入各节文件。 \\
    \texttt{latex/preamble.tex} & 宏包、页面格式、符号宏、占位宏与算法/代码环境。 \\
    \texttt{latex/refs.bib} & 参考文献数据库。 \\
    \texttt{latex/sections/} & 摘要、问题重述、预处理、三问、结论与附录。 \\
    \texttt{src/} & Python 脚本：数据读取、预处理、建模、绘图。 \\
    \texttt{img/} & 所有图片的输出目录，供 LaTeX 以 \texttt{../img/} 引用。 \\
    \texttt{附件/} & 赛题原始数据，只读；不在其中生成任何文件。 \\
    \bottomrule
  \end{tabularx}
\end{table}

\subsection{运行环境与依赖}

表~\ref{tab:appendix-env} 汇总本文实际使用的运行环境与依赖版本；从原始数据到论文图表的流程可在该环境中复跑。

\begin{table}[htbp]
  \centering
  \caption{运行环境与依赖}
  \label{tab:appendix-env}
  \small
  \begin{tabularx}{\textwidth}{@{}l X l@{}}
    \toprule
    项目 & 说明 & 版本 \\
    \midrule
    操作系统 & Ubuntu 22.04.5 LTS（内核 6.8.0-138-generic） & — \\
    Python & 解释器版本 & 3.10.12 \\
    科学计算 & \texttt{numpy}、\texttt{scipy}、\texttt{pandas} & 2.2.6 / 1.15.3 / 2.2.3 \\
    机器学习 & \texttt{scikit-learn}、\texttt{lightgbm}、\texttt{statsmodels} & 1.5.2 / 4.5.0 / 0.14.4 \\
    图与网络 & \texttt{matplotlib}、\texttt{networkx} & 3.9.2 / 3.3 \\
    表格读取 & \texttt{openpyxl} & 3.1.5 \\
    论文编译 & TinyTeX（TeX Live 2026）+ \texttt{xelatex} + \texttt{bibtex} & 2026 \\
    \bottomrule
  \end{tabularx}
\end{table}

本文的数据处理、建模与绘图使用 NumPy\citep{harris2020array}、pandas、scikit-learn\citep{pedregosa2011scikit}、NetworkX\citep{hagberg2008exploring} 与 Matplotlib\citep{hunter2007matplotlib}；具体版本以实际环境为准。

\subsection{复现步骤}

复现流程遵循“原始数据只读—脚本生成图片—LaTeX 引用图片”的顺序：

\begin{enumerate}
  \item 在项目根目录准备 Python 环境并安装表~\ref{tab:appendix-env} 中的依赖；
  \item 运行 \texttt{src/} 下的数据读取与预处理脚本，生成对齐后的数据与缺失掩码；
  \item 依次运行问题一至问题三的建模脚本，把结果表与图片写入约定目录；
  \item 进入 \texttt{latex/} 目录，执行 \texttt{latexmk -xelatex main.tex}；
  \item 检查生成的 PDF 中是否仍有红色占位文字，图片是否来自 \texttt{../img/}。
\end{enumerate}

命令示例见代码清单~\ref{lst:pipeline}。具体脚本名与参数以 \texttt{src/} 中的实际文件为准。

\begin{lstlisting}[style=shellcode,caption={复现命令},label={lst:pipeline}]
# Run from the project root
MPLBACKEND=Agg PYTHONPATH=.pylibs:src python3 src/01_profile.py
MPLBACKEND=Agg PYTHONPATH=.pylibs:src python3 src/00_figures.py
MPLBACKEND=Agg PYTHONPATH=.pylibs:src python3 src/02_external.py
MPLBACKEND=Agg PYTHONPATH=.pylibs:src python3 src/q1.py
MPLBACKEND=Agg PYTHONPATH=.pylibs:src python3 src/q2.py
MPLBACKEND=Agg PYTHONPATH=.pylibs:src python3 src/q3.py
# Build the paper (TinyTeX installed at ~/.TinyTeX)
export PATH="$HOME/.TinyTeX/bin/x86_64-linux:$PATH"
cd latex
latexmk -xelatex main.tex
\end{lstlisting}

\subsection{主要超参数与复现配置}

表~\ref{tab:appendix-config} 汇总影响结果的主要配置；论文中的数字均由同一配置复跑得到。

\begin{table}[htbp]
  \centering
  \caption{主要超参数与复现配置}
  \label{tab:appendix-config}
  \small
  \begin{tabularx}{\textwidth}{@{}l X X@{}}
    \toprule
    模块 & 配置项 & 取值 \\
    \midrule
    全局 & 随机种子、交叉验证折数 & 42；5 折 $\times$ 3 次重复 \\
    预处理 & 插补方法、MCAR/MAR 网格 & 折内中位数插补；MCAR 5/10/20/30\%；MAR 由 ethanol 分位数驱动 \\
    问题一 & 基学习器与校准 & RidgeCV、PLS 5 成分、ET/RF 500 棵、LGBM 400 棵（lr 0.03、叶子 15）；保序 + q01/q99 min-max \\
    问题二 & 图与稳定性 & $K=10$；$\alpha=0.3$；$B=500$；入选频率门槛 0.8；机制先验 $\lambda=0.5$ \\
    问题三 & 生成与缺陷约束 & 70 个校准化合物；log1p；Ledoit--Wolf；每风格 300 条；style cap $\max(\tau_j,q_{g,j})$ \\
    评价 & 指标与实现 & 宏平均 $R^2$/RMSE/MAE/Spearman；QWK；MMD；RF 风格分类器 \\
    \bottomrule
  \end{tabularx}
\end{table}

\subsection{核心代码说明}

表~\ref{tab:appendix-code} 给出本文代码模块的职责与主要输出；每个模块均可单独运行。

\begin{table}[htbp]
  \centering
  \caption{核心代码模块}
  \label{tab:appendix-code}
  \small
  \begin{tabularx}{\textwidth}{@{}l X X@{}}
    \toprule
    模块 & 功能 & 主要输出 \\
    \midrule
    数据读取与对齐 & \texttt{src/data\_utils.py}、\texttt{src/01\_profile.py}、\texttt{src/02\_external.py} & 对齐数据、变换表、EDA 图 \\
    问题一 & \texttt{src/q1.py}：岭回归/PLS/ET/RF/LGBM 等权融合、保序校准、缺失模拟 & 评分表、指标表、\texttt{img/q1\_*} \\
    问题二 & \texttt{src/q2.py}：共表达图、五方法融合、500 次 bootstrap、机制先验 & Top-3 表、网络图、稳定性表 \\
    问题三 & \texttt{src/q3.py}：条件 copula、OT、style cap、感官画像 & 生成谱、雷达图、评价表 \\
    概念图 & \texttt{src/00\_figures.py}：路线图、预处理流程图、Q1 框架图 & \texttt{img/fig\_*} \\
    \bottomrule
  \end{tabularx}
\end{table}

\subsection{可复现性检查清单}

提交前逐项核对：（1）随机种子是否固定并在附录记录；（2）数据划分是否在建模前固定，预处理是否只在训练折内拟合；（3）论文中的每个数字是否能由代码输出直接对应；（4）图片是否由脚本生成而非手工拼接；（5）\texttt{latex/} 中是否没有散落的中间文件；（6）是否没有把原始数据、模型权重或日志提交进仓库。


\end{document}
